\documentclass[usenatbib,twoside,twocolumn]{mnras}

% Paper packages
\usepackage{graphicx}
\usepackage{amsmath}
\usepackage{amssymb}
\usepackage{booktabs}

% Define astronomical units
\newcommand{\msun}{\ensuremath{M_{\odot}}}
\newcommand{\pc}{\ensuremath{\rm pc}}
\newcommand{\kms}{\ensuremath{\rm km\,s^{-1}}}

\title{Core Spacing in Herschel Gould Belt Survey Filaments: a Convention-Limited Sub-Classical Scale and its Interpretation}

\author[G.~J.~White]{G.~J.~White$^{1,2}$\\
$^{1}$School of Physical Sciences, The Open University, Walton Hall, Milton Keynes, MK7 6AA, UK\\
$^{2}$Rutherford Appleton Laboratory, Harwell Science and Innovation Campus, Didcot, OX11 0QX, UK}

\begin{document}

\maketitle

\begin{abstract}
\textit{Herschel} finds dense cores spaced along filaments well below the classical Inutsuka \& Miyama wavelength of an isothermal equilibrium cylinder. We re-analyse HGBS core spacings in four clouds (Gaia DR3 distances) with an Athena++ MHD campaign, reaching a deliberately limited conclusion: the spacing is sub-classical in \emph{sign} under every convention, but its magnitude is dominated by measurement convention, not physics or statistical error. The nearest-neighbour spacing is the appropriate observable; the common pairwise median instead measures the region extent $L/3$. Our single preferred value is $\lambda/W\approx1.4$ (region-Plummer NN), within a systematic band $\approx0.8$--$1.9$ across conventions; the robust result is the sub-classical \emph{sign}, a factor $\approx3$--$4$ below the classical $\approx5$--$6\times$FWHM, not the precise value. The cores are random-to-clustered, not periodic, so the observable is a core \emph{line density}, not a resonant wavelength. In simulations, contracting near-critical filaments fragment sub-classically --- a resolution-stable \emph{upper bound} $\lambda/W\lesssim1.0$--$1.3$, the spectrum not turning over --- while the identical pipeline recovers the classical wavelength for a numerically relaxed, force-balanced equilibrium, indicating the deficit is physical; being a bound, this fixes the sign, not sufficiency. Of the magnetic geometries, an \emph{ideal-MHD} perpendicular field collapses monolithically without axial fragmentation (short low-resolution beading is numerical), while a helical field --- the only one that could shorten $\lambda$ --- is seeded, not force-balanced, leaving that mechanism untested. We do not identify the cause: current data cannot separate single-filament fragmentation from fibre substructure, two readings that may not be physically distinct. Our primary results are methodological.
\end{abstract}


\begin{keywords}
stars: formation -- ISM: structure -- ISM: filaments -- MHD -- turbulence
\end{keywords}

\section{INTRODUCTION}

The \textit{Herschel} Space Observatory established that filaments are the fundamental structures within which star formation occurs in molecular clouds \citep{Andre2010}. Two key observational results emerged: (1) a characteristic filament width of $\sim$0.1 pc across diverse environments \citep{Arzoumanian2011}, and (2) a critical mass-per-unit-length of $M_{\rm line,crit} \approx 16~M_\odot$/pc \citep{Ostriker1964}.

Classical fragmentation theory \citep{Inutsuka1992} predicts that isothermal gas cylinders fragment at a wavelength of $\lambda_{\rm m}\approx5$--$6\times$ the filament FWHM (\citealt{Fischera2012}; equivalently $4\times$ the flat inner diameter, \citealt{Inutsuka1992}); for the characteristic HGBS width of $0.10$\,pc this is a core spacing of $\sim0.5$\,pc. Our primary \emph{physical} estimate of the nearest-neighbour spacing is $\lambda/W\approx1.4$, using directly-fitted region-specific Plummer widths (Section~\ref{sec:widths}); the fixed-width convention gives $\lambda/W\approx2.1$ (raw; $1.9$ after the periodic injection--recovery correction), retained only for comparability with prior HGBS work (Table~\ref{tab:headline} is the decoder for all correction states; both are sub-classical by a factor $\approx3$--$4$ relative to the $\lambda_{\rm m}\approx5$--$6\times$FWHM prediction)\footnote{Throughout, $W\equiv W_{\rm fil}$ is the observed HGBS FWHM and the reference width for all observation--theory comparisons; the fixed-width convention sets $W\approx0.10$\,pc, the region-Plummer convention $W=0.11$--$0.16$\,pc, and the simulation core half-width $W_{\rm core}$ relates to $W_{\rm fil}$ by the normalisation T1 (Section~\ref{sec:width_correction}).}. We adopt Gaia DR3-based distances \citep{Zhang2023} and the nearest-neighbour (NN) statistic, the estimator that most directly measures the fragmentation wavelength, and the one already used in the HGBS catalogue papers \citep{Konyves2015,Konyves2020} (Section~\ref{sec:statistics}).

This apparent deficit relative to the classical $\lambda_{\rm m}\approx5$--$6\times$FWHM has been attributed to additional physics. We note that \citet{Arzoumanian2019} established the $\sim$0.1\,pc \emph{width}, not a spacing; the sub-classical \emph{spacing} deficit and its usual fibre/magnetic reading are from the fragmentation literature (e.g.\ \citealt{Fischera2012,Zhang2023cal}). Two mechanisms are usually invoked: \emph{hierarchical fragmentation}, in which the filaments are bundles of velocity-coherent fibres \citep{Hacar2013,Hacar2018,Yang2024} so that, if each fibre fragments near the classical scale, measuring across the bundle compresses the apparent filament-level spacing; and \emph{magnetic tension}, in which a magnetic field alters the most unstable wavelength \citep{Nakamura1993}. (We note the field geometry carefully. For an axial fragmentation mode, a \emph{longitudinal} field is compressed \emph{along} its own direction and exerts essentially no restoring force, leaving $\lambda_{\rm max}$ almost unchanged \citep[][consistent with our $\theta=0^\circ$ runs being nearly $\beta$-independent]{Nagasawa1987,Nakamura1993}; a static \emph{perpendicular} field is amplified by the axial compression and \emph{lengthens}/stabilises the mode; only a \emph{helical}/toroidal field, whose pinch destabilises the sausage mode, shortens it \citep{Fiege2000}. ``Magnetic tension'' as a route to a shorter spacing is therefore a specifically helical effect; we test a \emph{seeded} helical field (Section~\ref{sec:optional_tests}) and find it narrows the width by a radial pinch rather than shortening $\lambda_{\rm max}$, but the \emph{force-balanced} helical equilibrium in which the \citet{Fiege2000} pinch mode operates is left to future work.) These are natural readings of the discrepancy within the equilibrium fragmentation theory that was the available comparison; what we suggest is not that they are wrong, but that an alternative comparison for a still-collapsing filament is the time-dependent one, which reduces the deficit without invoking them, though (as we argue) it does not by itself select among the readings.

From a large campaign of self-gravitating MHD simulations we measure the fragmentation wavelength of a dynamically collapsing filament from its growth-rate spectrum, and two caveats must be stated at the outset. First, the measured wavelength is an order-$\lambda_{\rm J}$ \emph{upper bound}, not a located peak: the growth-rate spectrum rises to a plateau at the short-wavelength edge of the resolution-converged band ($\lambda\gtrsim1\,\lambda_{\rm J}$) and does not turn over there, so $\lambda_{\rm max}\lesssim1\,\lambda_{\rm J}$ ($\lambda_{\rm max}/W_{\rm FWHM}\approx2$ versus the classical $\approx5$--$6$), and the true dynamical scale may be shorter (Appendix~\ref{sec:val_appendix}). Second, the like-for-like comparison must use matched widths: in a common beam-convolved Plummer convention the simulated value is $\lambda/W\approx1.0$--$1.3$, which lies \emph{between} the FilFinder ($0.4$--$0.8$) and DisPerSE ($1.9$, periodic-corrected; Table~\ref{tab:headline}) observed values and coincides with the directly-measured (raw) region-Plummer widths ($1.2$--$1.4$); the ``$\lambda_{\rm max}/W\approx2\approx1.9$'' agreement compares a Gaussian to a Plummer width and we do not rely on it. A filament that is still contracting is expected to fragment below the marginally-stable hydrostatic mode \citep{Clarke2016,Clarke2017}; our contribution is to reproduce that sub-hydrostatic behaviour numerically in Athena++ for the longitudinal near-critical case, and to connect it to a homogeneous Gaia-DR3 re-analysis of HGBS spacings. This is a numerical demonstration of the generic sub-classical fragmentation expected of any contracting, non-equilibrium filament, of which the accretion-driven \citet{Clarke2016,Clarke2017} mechanism is one realisation we do not specifically isolate. We solve the classical equilibrium dispersion relation only as the reference benchmark (Appendix~\ref{sec:val_appendix}); our contracting-filament wavelength is measured numerically, and the initial profile we adopt is not the exact isothermal equilibrium (Section~\ref{sec:limitations}), so we cannot fully separate the contracting-filament effect from an initial-condition contribution. The observations are therefore \emph{consistent with} single-filament dynamical fragmentation but do not \emph{require} it, nor discriminate it from hierarchical fibres or magnetic tension; which reading, if any, is favoured depends on the segmentation convention (Section~\ref{sec:error_budget}).

We present a full-coverage nearest-neighbour core-spacing analysis with Gaia DR3 distances \citep{Zhang2023} for the four robust HGBS regions (Orion~B, Aquila, Perseus, Taurus; the eight-region sample is retained only for the legacy pairwise-median statistic), followed by the simulation campaign, its validation and the interpretation. We note throughout that the large-scale filament--field misalignment from Planck polarimetry \citep{Planck2016} measures the projected orientation, which need not equal a filament's internal field geometry (Section~\ref{sec:discussion}).

\begin{table}
\caption{Principal symbols.}
\label{tab:symbols}
\footnotesize
\setlength{\tabcolsep}{4pt}
\begin{tabular}{lp{6.1cm}}
\toprule
Symbol & Meaning \\
\midrule
$\lambda$ & core-to-core spacing (an observed line-density scale, not necessarily a resonant wavelength; cf.\ Section~\ref{sec:clustering}) \\
$W_{\rm fil}$ & observed filament FWHM (beam-deconvolved Plummer; the reference width) \\
$W_{\rm core}$ & simulation initial Gaussian half-width $\sigma$ ($0.3\,\lambda_{\rm J}$) \\
$W_{\rm form}$ & FWHM of the evolved, projected simulation profile \\
$T_1\!\equiv\!\eta_W$ & $W_{\rm form}/W_{\rm fil}$, the simulation-to-observation width normalisation (forward-model band $0.59$--$0.78$; adopted $0.61$) \\
$\lambda_{\rm J}$ & background Jeans length ($=1$ in code units) \\
$\lambda_{\rm max}$ & dominant fragmentation wavelength (growth-rate peak) \\
$\lambda_{\rm IM92}$ & classical equilibrium-cylinder wavelength $\approx5$--$6\,W_{\rm FWHM}$ ($=4\times$ flat diameter) \\
$\Gamma(\lambda)$ & linear growth rate of axial mode of wavelength $\lambda$ \\
$f$ & thermal line-mass fraction $M_{\rm line}/M_{\rm line,crit}$ \\
$\beta$ & plasma beta $8\pi P/B^2$ \\
$\mathcal{M}$ & amplitude label of the seed perturbation ($\delta v=\mathcal{M}c_s\times10^{-4}$; a linear seed, \emph{not} a turbulent Mach number --- driven transonic turbulence is tested separately in Appendix~\ref{sec:supplementary_campaigns}) \\
$\mu_\Phi/\mu_{\Phi,\rm crit}$ & normalised mass-to-flux ratio \\
$\theta$ & angle between field and filament axis \\
$N_J$ & cells per local Jeans length (Truelove criterion) \\
$C$ & density contrast $\rho_{\rm max}/\langle\rho\rangle$ \\
\bottomrule
\end{tabular}
\end{table}

\section{OBSERVATIONAL ANALYSIS}

\subsection{Complete HGBS Sample}

Our sample comprises 8 high-confidence HGBS regions (Table~\ref{tab:sample}). The original catalogues \citep{Konyves2015, Konyves2020, Ladjelate2020, Pezzuto2021} adopted distances of $130$--$400$\,pc from the best estimates available in 2010--2015 (Table~\ref{tab:literature} lists the pre-Gaia distance and its source per region; note that \citealt{Konyves2020} already adopted $\sim400$\,pc for Orion~B); Gaia DR3 parallaxes \citep{Zhang2023} have since enabled substantial revisions, largest for the more distant regions, whose pre-Gaia distances were the most uncertain (Gaia fractional parallax errors are in fact largest, not smallest, at greater distance).

We adopt the Gaia DR3 YSO-based distances from \citet{Zhang2023} (Table~4), with typical uncertainties of 3\%--5\%. The revisions are strongly region-dependent. The largest are for Serpens ($260 \to 458$ pc, $+76\%$) and Aquila ($260 \to 436$ pc, $+68\%$); by contrast Orion~B, the \emph{largest} region in the sample, is essentially unchanged (its pre-Gaia HGBS distance was already $\sim400$\,pc, \citealt{Konyves2020}, and Gaia DR3 gives $386$\,pc, a $-3.5\%$ change), and Taurus ($140 \to 135$ pc) and Ophiuchus ($130 \to 137$ pc) barely move. Since spacings scale linearly with distance, the \emph{core-count-weighted} mean spacing therefore rises only modestly ($\sim10\%$ for the robust regions), and is dominated by Aquila rather than being a uniform $+32\%$ shift. (This $\sim10\%$ is the core-count-weighted mean, dominated by the unchanged Orion~B; the \emph{unweighted} per-region NN characteristic shifts more, $\lambda/W\approx1.4\to1.9$, because there the large Aquila revision carries equal weight.) Throughout, $\lambda$ denotes the observed core-to-core \emph{nearest-neighbour spacing}; where we mean the theoretical dominant unstable \emph{wavelength} we write $\lambda_{\rm max}$ (simulations) or $\lambda_{\rm m}$ (classical). The observed spacing is a line-density statistic, not a resonance, and the two coincide only to order of magnitude.

The analysis includes all HGBS-identified cores (starless, prestellar and protostellar), extracted and classified within the consistent HGBS framework \citep{Konyves2015, Konyves2020}. Including unbound starless and protostellar sources may introduce some bias, but fragmentation produces a continuum of boundness states, and restricting the sample to prestellar cores would shrink several regions (notably TMC1, CRA and Serpens) below statistical usefulness. Core positions and properties are taken directly from the published catalogues.

Protostellar migration displacements of $0.01$--$0.05$\,pc \citep{Kirk2016, Mattern2018} are small against the $\sim0.2$--$0.3$\,pc spacing; we estimate the resulting NN bias at $\sim5$--$10\%$ and carry it as a systematic (Section~\ref{sec:error_budget}).

\begin{table*}
\caption{Complete HGBS Sample with Gaia DR3 Distances. Per-region spacings are pairwise-median values rounded to three significant figures; the core-count-weighted mean is $0.279$\,pc (full sample) and $0.290$\,pc (robust), the latter used only for the robust-vs-full comparison (Section~2.3).}
\label{tab:sample}
\begin{tabular}{lccccc}
\toprule
Region & Distance & Total & Pairwise ($L/3$) & Std Error & Status \\
       & (pc)   & Cores & (pc)   & (pc)      &  \\
\midrule
Orion B  & 386 & 1,844 & 0.313 & 0.047 & Robust \\
Aquila   & 436 & 749   & 0.346 & 0.047 & Robust \\
Perseus  & 300 & 816   & 0.248 & 0.040 & Robust \\
Taurus   & 135 & 536   & 0.198 & 0.040 & Robust \\
Ophiuchus& 137 & 513   & 0.206 & 0.053 & Limited \\
Serpens  & 458 & 194   & 0.331 & 0.097 & Limited \\
TMC1     & 135 & 178   & 0.195 & 0.056 & Limited \\
CRA      & 150 & 239   & 0.248 & 0.072 & Limited \\
\midrule
\textbf{Weighted Mean}$^\dagger$ &  & \textbf{5,069} & \textbf{0.279} & \textbf{0.009} &  \\
\bottomrule
\end{tabular}
\par\vspace{2pt}
\footnotesize $^\dagger$The $0.009$\,pc figure is a formal core-count-weighted standard error (treating cores as independent); it is \emph{not} the relevant uncertainty, since the four clouds are the independent units (Section~\ref{sec:error_budget}). It is retained only as a within-region formal quantity and plays no role in the comparison with theory.
\end{table*}

\subsection{Measurement Methodology}
\label{sec:methodology}

The analysis pipeline, from raw HGBS maps to the final characteristic $\lambda/W$, is summarised in Fig.~\ref{fig:pipeline}. Filament skeletons were identified using DisPerSE \citep{Sousbie2011} following standard HGBS methodology \citep{Arzoumanian2011}, with a persistence threshold of $3\sigma$ above the background noise and manual editing restricted to removing obvious artefacts (e.g.\ spurs connecting unrelated structures). This manual editing was deliberately limited to the excision of clear artefacts and did not alter genuine filament topology; its effect on the measured spacing is small, $<5\%$, as the junction-retention and association-threshold sensitivity tests below bracket the same class of skeleton perturbation. We nonetheless flag it as an \emph{unquantified} systematic: it has not been bounded by independent re-editing or by blind tests, and a blind, multi-editor reproducibility test of the skeletonisation is a recommended future check. Because the skeletons and core positions derive from angular positions on the sky, which are distance-independent, the Gaia DR3 revisions require no re-extraction: physical scales follow by rescaling the angular positions, leaving core associations and skeleton topology unchanged. Cores at filament junctions are retained, following standard HGBS practice; junctions are physically relevant sites where filaments converge, and excluding them would introduce selection bias about what counts as a filament segment. Cores were associated with filaments within a perpendicular threshold of $2\times$ the characteristic width ($0.2$\,pc, applied as a fixed physical distance). The measured spacing is insensitive to these choices: excluding junctions changes it by $<5\%$, and varying the association threshold between $1\times$ and $1.5\times$ the width changes it by $<3\%$, both well within the statistical uncertainty. (This $2\times$-width, $0.2$\,pc threshold refers to the pairwise-median association; the primary along-skeleton NN pipeline instead associates cores within a filament half-width, $\sim0.06$\,pc, of a spine, Section~\ref{sec:injection_recovery}.)

Core-to-core spacings were also measured using the pairwise median statistic, which Section~\ref{sec:statistics} shows converges to $L/3$ rather than the fragmentation wavelength for $N\gtrsim5$; we retain it only as a comparability statistic with prior HGBS work.

Because physical spacing is angular separation times distance, the $\lambda\propto d$ scaling is tautological and we do not present it as corroboration; the relevant check is that the \emph{residuals} about that scaling do not correlate with the absolute revision fraction $|\Delta D/D_{\rm original}|$ ($r=0.18$, $p=0.68$), so the revision does not preferentially bias the most-revised clouds. Two caveats remain, and we flag them as limitations rather than resolve them: the YSO-based distances \citep{Zhang2023} are assumed to apply to the dense gas measured here, and the Aquila/Serpens Rift is known to have substantial line-of-sight depth, so assigning a single distance may blend populations and inflate the inferred spacing. As a floor, retaining the pre-Gaia distances lowers the robust-region characteristic to $\lambda/W\approx1.4$ (still sub-classical), so the sub-classical \emph{sign} does not depend on the distance revision, though its magnitude does.

\begin{figure}
\centering
\begin{minipage}{0.95\columnwidth}
\footnotesize
\begin{enumerate}\setlength{\itemsep}{2pt}
\item HGBS column-density maps ($250\,\mu$m reference) + Gaia DR3 distances
\item DisPerSE skeleton extraction ($3\sigma$ persistence threshold)
\item Morphological closing (8 iterations) $+$ re-skeletonisation (bridge fragmented segments; components $\geq40$\,px)
\item Core association: cores within $0.06$\,pc of a spine are assigned; others rejected
\item Graph ordering: cores ordered by double-BFS arclength along each component
\item Adjacent-core NN spacing (in pc, via Gaia DR3 distance)
\item Injection--recovery bias correction ($\div\,1.11$--$1.29$, region-dependent)
\item Characteristic $\lambda/W = $ bias-corrected median spacing $/\,W_{\rm fil}$
\end{enumerate}
\end{minipage}
\caption{Analysis pipeline from raw HGBS maps to the final characteristic $\lambda/W$. Each step's parameters are given in the text (Sections~\ref{sec:statistics}--\ref{sec:injection_recovery}).}
\label{fig:pipeline}
\end{figure}

\subsection{Results}

The four robust regions are the independent statistical units throughout this paper; individual cores and spacings within a cloud are not independent, and we propagate this nesting in the hierarchical bootstrap (Section~\ref{sec:error_budget}). With only four independent clouds we treat the cloud-to-cloud scatter ($\sigma\approx0.4$) as a descriptive range, not a formal confidence interval; we do not attach a bootstrap CI to an $N=4$ sample.

\begin{table}
\caption{This table is the convention decoder used throughout: read every $\lambda/W$ quoted in the text against it. The nearest-neighbour $\lambda/W$ under each width/skeleton convention and correction state, for the four robust HGBS regions (revised distances). Every combination is sub-classical (classical $\lambda_{\rm m}\approx5$--$6\times$FWHM). We adopt the directly-measured region-Plummer value ($\lambda/W\approx1.4$) as the primary physical headline; the fixed-width value ($\lambda/W\approx1.9$) is retained for comparability with prior HGBS work. The ``bias-corr.'' column applies the (periodic) injection--recovery correction; the data-matched \emph{clustered} correction (Section~\ref{sec:clustering}) is a further downward factor $\div1.64$, giving $\approx1.2$ (fixed width) and $\approx0.8$ (region-Plummer), so every convention and correction is sub-classical. The classical fastest-growing wavelength is $\lambda_{\rm m}\approx5$--$6\times$FWHM (\citealt{Fischera2012}; $=4\times$ the flat diameter, \citealt{Inutsuka1992}); the DisPerSE conventions are sub-classical by a factor $\approx3$--$4$ (FilFinder more) (Appendix~\ref{sec:val_appendix}).}
\label{tab:headline}
\footnotesize
\setlength{\tabcolsep}{4pt}
\begin{tabular}{p{3.6cm}ccc}
\toprule
Convention & raw & periodic & clustered \\
\midrule
Region Plummer, DisPerSE (\emph{primary}) & $1.4$ & $1.2$ & $0.8$ \\
Fixed $0.10$\,pc, DisPerSE & $2.1$ & $1.9$ & $1.2$ \\
FilFinder skeleton & $0.5$--$0.9$ & $0.4$--$0.8$ & $0.3$--$0.5$ \\
Pairwise median ($L/3$, not $\lambda$) & $2.8$ & --- & --- \\
\midrule
Matched simulation (upper limit) & \multicolumn{3}{c}{$\lambda/W\lesssim1.0$--$1.3$} \\
Classical $\lambda_{\rm m}$ (F\&M 2012) & \multicolumn{3}{c}{$5$--$6$ ($=4\times$ diam.)} \\
\bottomrule
\end{tabular}
\\[2pt]
\footnotesize
\emph{Notes} The primary physical result is the directly-measured region-Plummer NN $\lambda/W\approx1.4$ (plausible range $0.8$--$1.9$ across the conventions and injection--recovery corrections above; not a formal confidence interval; Section~\ref{sec:error_budget}). The fixed-width convention gives a \emph{higher} value ($\approx1.9$, or $\approx1.2$ under the data-matched clustered correction; Section~\ref{sec:clustering}) because it uses $0.10$\,pc rather than the larger measured widths; the FilFinder skeleton gives lower values still. The pairwise median converges to $0.293\,L\approx0.88\,(L/3)$ (the median of the triangular pairwise-distance law, whose \emph{mean} is $L/3$; Section~\ref{sec:statistics}) --- a region-extent scale, not a spacing --- and empirically overestimates the fragmentation wavelength by $\sim$45--50\% here, so it is not directly comparable to the classical prediction. Per-region values are given in Sections~\ref{sec:widths} and~\ref{sec:injection_recovery}. Given the convention systematic, we quote $\lambda/W$ to one decimal place; finer digits are not significant.
\end{table}

\textbf{Robust regions.} We adopt the four ``robust'' regions (Orion~B, Aquila, Perseus, Taurus; $N>500$ cores each, well-established distances) as the primary sample.

\textbf{Full sample result.} The core-count-weighted mean across all 8 HGBS regions ($0.279$\,pc) and the robust-only weighted mean ($0.290$\,pc) differ by $\approx3.9\%$, and a leave-one-out analysis confirms no single region dominates ($<6\%$ shifts). All spacings are quoted to three significant figures; where an exact recomputation from the rounded per-region inputs differs from the bootstrap-derived primary value it does so by $\lesssim1\%$, which does not affect any conclusion.

All weighted means use core-count weighting,
\begin{equation}
\bar{\lambda} = \frac{\sum_i N_i\,\lambda_i}{\sum_i N_i},
\label{eq:weighting}
\end{equation}
where $\lambda_i$ and $N_i$ are the per-region spacing and core count.

The Gaia DR3 distance revisions raise the core-count-weighted mean spacing modestly and region-dependently ($\sim10\%$ for the robust regions; large for Aquila/Serpens, negligible for the dominant Orion~B; Section~\ref{sec:distance_uncertainties}); a uniform $260$\,pc baseline would instead imply a $+32\%$ shift, which over-states the true, region-weighted revision. The absolute $\lambda/W$ result below uses the adopted Gaia distances and is unaffected by this. With only four independent clouds we do not attach a formal confidence level to the comparison: a nonparametric resampling of the four regions is degenerate and cannot estimate a population variance, so we quote the four per-region values and their min--max range rather than a bootstrap ``95\% CI'' (resampling the four pairwise-median values spans $0.261$--$0.298$\,pc, which we treat as descriptive only). The preferred NN measurement is presented in Section~\ref{sec:injection_recovery}.

\subsection{The dominant observational systematic: filament width}
\label{sec:widths}

The fixed $W_{\rm fil}=0.10$\,pc convention is the largest single lever on $\lambda/W$, comparable to the whole sub-classical discrepancy, so we measured region-specific widths directly. Along each skeleton spine we sampled the column-density profile perpendicular to the local tangent and fitted a beam-deconvolved Plummer function, the HGBS width definition \citep{Arzoumanian2019}, to hundreds of profiles per region ($18.2''$ beam; Table~\ref{tab:widths}). All four robust regions give a well-defined Plummer index $p\simeq2.0$ and a width $W_{\rm fil}=0.11$--$0.16$\,pc, agreeing with the published HGBS values \citep{Arzoumanian2019} and slightly broader than the canonical $0.10$\,pc. The region-specific $\lambda/W$ is $1.1$--$1.8$ (median $\approx1.4$), so the region-specific convention gives a \emph{lower} value than the fixed-width headline ($\approx2.1$ (raw; Table~\ref{tab:headline})), not a higher one, because the true widths exceed $0.10$\,pc; every robust region is sub-classical under either convention. A Gaussian-crest width estimate returns a spuriously narrow Taurus width ($0.031$\,pc, formally driving $\lambda/W\to4$); the physically-motivated Plummer fit ($W_{\rm fil}=0.11$\,pc, $\lambda/W\approx1.1$) is the appropriate measure and makes Taurus the most sub-classical region rather than a classical one. The sub-classical result is thus not conditional on the width convention.

\begin{table}
\caption{Region-specific, beam-deconvolved filament widths from a direct Plummer fit to perpendicular column-density profiles along the skeletons (the HGBS width definition; \citealt{Arzoumanian2019}), and the resulting NN $\lambda/W$. $W_{\rm fil}$ is the beam-deconvolved Plummer FWHM ($18.2''$ beam), $p$ the Plummer index (a \emph{free} fit parameter, not fixed; it converges to $p=2.00$--$2.01$ in all four regions, consistent with the canonical Plummer-2 profile), $N$ the number of fitted profiles. All four regions give $p\simeq2$ and sub-classical $\lambda/W$.}
\label{tab:widths}
\footnotesize
\setlength{\tabcolsep}{3pt}
\begin{tabular}{p{1.7cm}ccccc}
\toprule
Region & $W_{\rm fil}$ (pc) & $p$ & $N$ & NN (pc) & $\lambda/W$ \\
\midrule
Taurus  & $0.110$ & $2.01$ & 295 & $0.124$ & $1.1$ \\
Orion B & $0.148$ & $2.00$ & 356 & $0.207$ & $1.4$ \\
Perseus & $0.147$ & $2.00$ & 235 & $0.263$ & $1.8$ \\
Aquila  & $0.161$ & $2.00$ & 446 & $0.210$ & $1.3$ \\
\midrule
median  & $0.148$ & $2.00$ & --- & --- & $1.4$ \\
\bottomrule
\end{tabular}
\par\vspace{2pt}
\footnotesize The measured widths ($0.11$--$0.16$\,pc) agree with the published HGBS values and exceed the canonical $0.10$\,pc, so region-specific Plummer widths give a \emph{lower} $\lambda/W$ ($\approx1.4$, raw) than the raw fixed-width headline ($\approx2.1$; Table~\ref{tab:headline}): all four regions are sub-classical under either convention. NN spacings are the full-coverage medians (Section~\ref{sec:injection_recovery}), measured on the same DisPerSE skeleton as the widths (its deposited crest for the width fit, its bridged form for the spacing); the injection--recovery bias correction lowers each $\lambda/W$ by a further $\sim$10--20\%. (Pipeline outputs retained to 2 s.f.; the convention systematic dominates.)
\end{table}

The fit is consistent across the sample: every region returns a Plummer index $p\simeq2.0$, the value found for HGBS filaments generally \citep{Arzoumanian2019}, and a width within $\sim50\%$ of $0.10$\,pc (inter-quartile ranges $\sim0.07$--$0.20$\,pc), so the sub-classical $\lambda/W$ is a property of all four robust regions, not of one anomalous cloud. The single earlier outlier, a spuriously narrow $0.031$\,pc Gaussian crest for Taurus, was a profile-selection effect of fitting the inner ridge rather than the broad Plummer wings (not a pixel-scale artefact: the $3''$ pixels subtend $0.002$\,pc at $135$\,pc); it does not survive the Plummer fit, which gives $0.11$\,pc and $\lambda/W\approx1.1$.

\subsection{Sample Classification and Distance Robustness}
\label{sec:distance_uncertainties}

We classify the eight regions (Table~\ref{tab:sample}) as \textbf{robust} (Orion~B, Aquila, Perseus, Taurus) or \textbf{limited} by three criteria: $N>500$ cores; a well-established distance; and a skeleton simple enough for unambiguous core-to-filament association. Limited regions fail at least one, Ophiuchus on skeleton complexity, the others on small samples, and are excluded from the primary NN measurement. The largest distance revisions (Serpens $+76\%$, Aquila $+68\%$) are independently corroborated: direct VLBI maser parallaxes validate Gaia DR3 to $<5\%$ where available \citep{Kounkel2018,Xu2021,OrtizLeon2018,OrtizLeon2017}, and for Serpens the 458\,pc value is supported by extinction-based estimates \citep[$440$--$460$\,pc;][]{Yan2022,Green2024} and by co-spatiality with Aquila in the Aquila Rift \citep{Zucker2020}. Serpens is a \emph{limited} region, excluded from the primary four, so its co-spatiality with Aquila in the Rift does not enter the four-cloud independence assumption; any residual line-of-sight blending would only \emph{reduce} an effective sample size we already conservatively take as four. The sub-classical conclusion is independent of these two regions: a leave-one-out analysis shifts the weighted mean by $<6\%$, reverting the limited regions to their original HGBS distances gives a pairwise-median $\lambda/W=2.68$ (sub-classical relative to $\lambda_{\rm m}\approx5$--$6\times$FWHM), and the distance-revision consistency check of Section~\ref{sec:statistics} confirms the $\lambda\propto d$ unit conversion. The characteristic is not driven by Aquila (the largest distance revision and fewest NN pairs of the robust four): dropping it leaves the region-Plummer median at $1.40$ (from $1.35$) and the fixed-width characteristic at $2.1$, both still sub-classical.


\subsection{Choice of spacing statistic: the $L/3$ property of the pairwise median}
\label{sec:statistics}

We adopt the nearest-neighbour (NN) spacing as the primary measure of the fragmentation wavelength, reporting the all-pairs median only for comparability with prior HGBS work. The two differ fundamentally because the pairwise median has an analytic $L/3$ convergence property. For $N$ points uniformly distributed on a filament of length $L$, a single pairwise distance $D=|X_i-X_j|$ has the triangular PDF $f(d)=2(L-d)/L^2$ ($d\in[0,L]$), with mean $\langle D\rangle=L/3$ and median $d^\star=L(1-1/\sqrt{2})\approx0.293\,L\approx0.88\,(L/3)$. The median of all $N(N-1)/2$ pairwise distances converges to $d^\star$ as $N\to\infty$, for a random or a perfectly periodic distribution alike, so on an idealised filament it measures the extent, not the fragmentation wavelength. Real catalogues, with finite sampling, branching and incomplete skeletons, deviate from this limit, so we treat $L/3$ as the statistic's limiting behaviour rather than an exact per-filament prediction; in practice it tracks $\sim L/3$ for any filament with $N\gtrsim5$ cores, and our injection--recovery test reproduces the $\approx0.88\,(L/3)$ asymptote (Section~\ref{sec:injection_recovery}). The finite-$N$ behaviour is a small positive bias: Monte-Carlo over $N=5$--$500$ gives median$/L\approx0.293\,(1+0.43/N)$, i.e.\ $\le2\%$ inflation for any filament with $N\gtrsim20$ cores, so the statistic tracks $0.293\,L$ to a few per cent for realistic $N$ and never approaches the fragmentation wavelength ($\sim30$--$100\times$ smaller). For \emph{branched} skeletons the along-skeleton pairwise distances are geodesic, so the statistic measures $\approx0.3$ times the connected component's arclength \emph{extent} (diameter) rather than its total length --- across linear, Y, star, comb and binary-tree topologies the median stays $0.89$--$1.25\times$(diameter$/3$); and for a population with a realistic length distribution the pooled median is pair-weighted toward the longest filaments and tracks the rms-length$/3$, overestimating any true wavelength by $17$--$35\times$. The pairwise median is thus an extent estimator under finite $N$, branching, and realistic length mixes alike (full tables in Appendix~\ref{sec:val_appendix} deposit / online material). The HGBS catalogue papers \citep{Konyves2015,Konyves2020} and fibre-resolved analyses \citep{Hacar2013,Hacar2018} likewise use nearest-neighbour separations. The NN statistic is unaffected by this bias and is beam-safe: the $18.2''$ HGBS beam corresponds to $0.012$--$0.039$\,pc across our robust regions, $\gtrsim6\times$ smaller than the measured NN spacing.

We reconstruct the along-skeleton core ordering for the four robust regions from the deposited DisPerSE skeleton maps, bridging fragmentation by morphological closing, and extend the NN measurement to all four at full coverage in Section~\ref{sec:injection_recovery}; the limited regions are excluded from the NN measurement owing to branching skeletons or small samples. Figure~\ref{fig:spacing} shows the per-region spacings with revised distances.

\begin{figure*}
\includegraphics[width=0.95\columnwidth]{figures/figure1_spacing_comparison.pdf}
\caption{Primary nearest-neighbour core spacing $\lambda/W_{\rm fil}$ per robust region, in the two DisPerSE width conventions: directly-measured region-Plummer widths (blue circles, our preferred; per-region raw $1.12,1.40,1.79,1.31$, median $1.4$, dashed) and the fixed $0.10$\,pc width (orange squares; $1.2,2.1,2.6,2.1$), both shown \emph{raw} (before the uniform injection--recovery correction, which lowers them; Section~\ref{sec:injection_recovery}). The red band is the classical fastest-growing wavelength $\lambda_{\rm m}\approx5$--$6\times$FWHM (\citealt{Fischera2012}; $=4\times$ the flat diameter, \citealt{Inutsuka1992}), and the dotted line the classical \emph{critical} wavelength $\lambda_{\rm cr}\approx2.6\times$FWHM; every region lies well below both. The grey band marks the independent FilFinder skeleton range ($\lambda/W\approx0.4$--$0.9$). The deficit is sub-classical in sign under every convention; its magnitude is dominated by the width/skeleton choice, not statistical error (Section~\ref{sec:error_budget}).}
\label{fig:spacing}
\end{figure*}

\subsection{Injection--recovery validation and model dependence of the correction}
\label{sec:injection_recovery}

To test whether the nearest-neighbour (NN) statistic recovers known spacings under realistic HGBS conditions, we placed synthetic cores with known periodic spacings (0.15--0.35\,pc) on HGBS-property-matched DisPerSE skeletons for all 8 regions and applied the NN statistic with identical methodology. Across 40 tests the raw measurement recovers the input with a small \emph{intrinsic-estimator} bias of $3.05 \pm 0.69$\% on clean skeletons. The larger correction actually applied to the headline is the \emph{full-coverage pipeline} bias ($\div1.11$--$1.29$, from junction/branch topology on the real bridged skeletons; Section~\ref{sec:injection_recovery}), which moves raw $2.1\to1.9$. Because that test injects a periodic pattern whereas real cores may be clustered, we also ran clustered injection--recovery (Poisson inflation $1.03$--$2.47$, mean $1.64$; lognormal $0.98$--$1.63$, mean $1.22$), and --- crucially --- we determined \emph{directly from the data} which regime applies (Section~\ref{sec:clustering}). The observed adjacent-core spacings are random-to-clustered, not periodic (coefficient of variation of order unity, $0.65$--$1.36$; a regular ``periodic$+$jitter'' distribution is rejected at $p\lesssim10^{-15}$ in every region). It is worth separating two conceptually distinct adjustments here. The \emph{topological} pipeline bias ($\div1.11$--$1.29$, from junction/branch ordering on the real bridged skeletons) is a genuine, region-measured correction and is applied. The \emph{point-process model factor} (periodic $\div1.20$ vs clustered $\div1.64$) is not a bias on a line-density observable at all: if the cores are genuinely random-to-clustered, the measured NN median already \emph{is} the physical quantity, and ``correcting'' it toward an injected periodic template the data reject would be circular. We therefore do not fold the model factor into the headline: we carry the \emph{raw} value as primary (Section~\ref{sec:injection_recovery_model}) and quote the clustered branch only as the matching lower bound. For completeness, were one to apply the clustered factor ($\div1.64$, the best match to the observed CoV$\approx1$) the raw apparent $\lambda/W\approx2.0$ would map to $\approx1.2$ (range $0.8$--$1.7$ across the clustered/lognormal family), whereas the periodic factor gives $\approx1.9$ (assuming a regularity the data reject). The sub-classical \emph{sign} holds under every choice.

Repeating the injection--recovery for the pairwise-median statistic confirms that it does not recover $\lambda_{\rm true}$: it returns a value $\sim$$20\times$ larger, tracking $L/3$ ($\langle\lambda_{\rm PW}\rangle/\lambda_{\rm true}\approx L/(3\lambda_{\rm true})$) and asymptoting to $\approx0.88\,(L/3)$, while the NN statistic recovers $\lambda_{\rm true}$ to $1.8$--$3\%$ on the identical skeletons. Computing the pairwise median over \emph{all} cores in each robust region (without per-filament segmentation) returns $\sim L_{\rm region}/3$ ($12.5$, $7.5$, $8.0$, $2.5$\,pc for Orion~B, Aquila, Perseus, Taurus), the region extent, not the $\sim0.2$--$0.35$\,pc per-filament values, confirming directly that the statistic measures segmentation scale rather than the fragmentation wavelength.

\textbf{Full-coverage along-skeleton NN for the robust regions (fixed-width along-skeleton NN)}. We extend the along-skeleton NN measurement to all four robust regions at full coverage using the dense deposited DisPerSE skeletons, which (unlike the sparse FilFinder spines) pass near essentially all cores. These skeletons are fragmented; we bridge the fragmentation by morphological closing followed by 1-pixel re-skeletonization, order cores along each reconstructed filament by graph arclength, and associate cores lying within a filament half-width ($\sim$0.06\,pc) of a spine. This yields $N_{\rm pairs} = 461,\,614,\,563,\,168$ adjacent-core spacings for Taurus, Orion~B, Perseus and Aquila respectively. The median NN spacings are $0.124,\,0.207,\,0.263,\,0.210$\,pc, i.e.\ $\lambda/W = 1.2,\,2.1,\,2.6,\,2.1$, with a robust-region characteristic of $\lambda/W \approx 2.1$. The result is robust to the bridging length: $\lambda/W$ for Orion~B stays in the range $1.9$--$2.7$ as the closing radius spans a factor $>2$, and is insensitive to the association radius. Injecting synthetic cores at known spacings onto the real bridged skeletons (with their genuine junction/branch topology) and recovering them with the identical pipeline shows a quantified overestimate of $\sim$$15$--$20$\% at the observed scale (recovery ratio $1.11$--$1.29$ across regions; larger at smaller injected spacings), arising from the longest-path ordering at junctions. The injection--recovery correction is a downward systematic whose magnitude is uncertain because the cores are non-periodic (Section~\ref{sec:clustering}): the periodic pipeline factor ($\div1.11$--$1.29$) gives $\lambda/W\approx1.0,\,1.9,\,2.1,\,1.8$ (Taurus, Orion~B, Perseus, Aquila; robust-region characteristic $\approx1.9$), while a clustered factor ($\div1.64$) would lower this to $\approx1.2$. We apply this correction \emph{uniformly} to every width convention rather than to one alone, so the fixed-width characteristic is $\lambda/W\approx1.2$--$1.9$ and the region-Plummer characteristic (Section~\ref{sec:widths}) is $\lambda/W\approx0.8$--$1.4$; we quote the directly-measured (raw) region-Plummer value $\lambda/W\approx1.4$ as the primary figure and carry the correction as a stated uniform systematic. We flag its consequence for the theory comparison plainly: under the raw or periodic values ($\lambda/W\approx1.2$--$1.4$) single-filament fragmentation (simulated $\approx1.0$--$1.3$) suffices, whereas under the data-favoured clustered correction the observed value falls to $\approx0.8$, below the simulated scale, where single-filament fragmentation would be insufficient and additional compression is required --- so the sign of the observation--simulation offset itself depends on the (uncertain) correction, reinforcing the under-determination conclusion. \textbf{The full-coverage NN spacing ($\lambda/W\approx1.2$--$2.1$ across corrections) is sub-classical --- a factor $\approx3$--$4$ below (DisPerSE conventions; more under FilFinder) the classical $\lambda_{\rm m}\approx5$--$6\times$FWHM --- and the pairwise-median value $\lambda/W\approx2.8$ overestimates it by $\sim$45--50\% through the $L/3$ extent-bias.} Taurus is the most sub-classical region, not an outlier.

As an independent cross-check, running the same along-skeleton BFS-arclength NN pipeline on independently extracted FilFinder spine maps (all four robust regions) gives $\lambda/W\approx0.4/0.6/0.7/0.8$ for Taurus, Orion~B, Perseus, Aquila respectively, systematically $2$--$3\times$ smaller than the DisPerSE values ($1.0/1.9/2.1/1.8$) because the FilFinder mask threshold produces a denser, more branched network. Both methods agree the characteristic spacing is far sub-classical ($\lambda/W\ll4$) in every region; FilFinder if anything strengthens the conclusion.

\smallskip\noindent\textbf{Model dependence of the correction.}\label{sec:injection_recovery_model}
The downward injection--recovery correction is not a determined number but is \emph{conditional} on the assumed underlying point-process model: a periodic injection returns a mild factor ($\div1.1$--$1.3$), whereas a clustered (Poisson/lognormal) injection returns a larger one (up to $\div1.6$). There is, moreover, an unavoidable circularity in selecting between them --- we observe that the cores are clustered (Section~\ref{sec:clustering}), adopt a clustered injection model on that basis, and thereby infer a smaller intrinsic spacing --- which we state explicitly rather than conceal. We therefore carry the \emph{raw} (uncorrected) spacing as the primary quantity throughout, and treat the injection--recovery factor as a stated, model-dependent systematic spanning $\div1.1$ (periodic) to $\div1.6$ (clustered), not as a single corrected value. We tested this directly. For each region we derived a correction matched to its own measured clustering (coefficient of variation of the along-skeleton core gaps, CoV$=1.36,0.83,0.96,0.65$ for Taurus, Orion~B, Perseus, Aquila) by injecting a gamma renewal process at that CoV and recovering it through a deblending pipeline anchored to the periodic ($\div1.20$) and Poisson ($\div1.64$) endpoints. The region-matched factors are $\div1.86,1.53,1.61,1.40$, giving corrected $\lambda/W\approx0.7,1.4,1.6,1.5$ (fixed $W=0.10$\,pc). As expected the pooled factor slightly over-corrects the more regular regions (Aquila, Orion~B) and under-corrects the clustered Taurus, but the effect on the robust-region headline is small (mean $\lambda/W$ $1.2\to1.3$): the sub-classical result holds in every region under its own matched correction, and $\lambda/W\approx1.9$ --- which requires the rejected periodic model --- remains excluded. We therefore retain the pooled value as the headline while noting it is a population average of a CoV spanning $0.65$--$1.36$.

\subsection{Regularity of the core spacing: periodic, random, or clustered?}
\label{sec:clustering}
Classical gravitational fragmentation predicts a \emph{quasi-periodic} chain of cores, whereas fibre projection or a turbulent/accretion-driven cascade predicts a broad, random-to-clustered distribution with no single characteristic comb. We tested this directly. Ordering the associated cores by arclength along each filament spine and measuring the adjacent-core gaps, we find a coefficient of variation $\mathrm{CoV}=\sigma/\mathrm{mean}$ of order unity in every robust region (Taurus $1.36$, Orion~B $0.83$, Perseus $0.96$, Aquila $0.65$; pooled $1.16$), with gap histograms that closely follow an exponential (Poisson) law (Fig.~\ref{fig:clustering}). Three of the four regions are Poisson-to-clustered ($\mathrm{CoV}\gtrsim0.8$); Aquila alone ($\mathrm{CoV}=0.65$) is mildly \emph{sub}-Poisson, i.e.\ weakly more regular than random, though it also has the fewest pairs ($168$) and lowest skeleton coverage ($24\%$), so we do not over-read it. A Kolmogorov--Smirnov test rejects a regular ``periodic$+$jitter'' distribution at $p\lesssim10^{-15}$ in all four regions --- a rejection that is robust to the estimated-parameter issue below, since fitting parameters only makes a null \emph{harder} to reject. The exponential model is an acceptable description in three regions and only marginally so in the fourth; we quote these goodness-of-fit $p$-values as \emph{indicative} only, because the exponential rate is estimated from the same data, which invalidates the standard KS null distribution (a Lilliefors or parametric-bootstrap calibration would be required for exact values). The load-bearing statistical statement is the \emph{rejection of periodicity}, which does not depend on this. A mild deficit of the smallest separations (1-D Clark--Evans ratio $R>1$) reflects a finite-core/deblending exclusion at the $\sim$beam scale, not quasi-periodicity (which additionally requires $\mathrm{CoV}\!\to\!0$). Quasi-periodic fragmentation is therefore \emph{not} strongly imprinted: the spacings are consistent with a random process bearing a mild hard-core exclusion and a weakly clustered tail. This is the theory-side statement of \citet{Clarke2020} --- fragmentation through sub-filaments leaves no single characteristic length --- seen directly in the data. It also means that \emph{if} an injection--recovery correction is applied (Section~\ref{sec:injection_recovery}), the clustered branch, not the periodic one, is the relevant one, giving a corrected lower bound $\lambda/W\approx1.2$; we nonetheless carry the raw value as the primary figure and treat the correction as a stated systematic (Section~\ref{sec:injection_recovery_model}). A shifted-exponential (minimum-separation) model fits the gaps better than either a pure exponential or a peaked distribution in every region (by AIC), with an inferred hard core of $0.011$--$0.044$\,pc $\approx$ the beam/core diameter. One statistical consequence of this near-exponential form deserves flagging: for an exponential gap distribution the \emph{median} gap is $\ln2\approx0.69$ times the \emph{mean} gap, and it is the mean gap that equals the inverse line density $1/\rho_{\rm lin}$ we argue is the physical observable. Our headline uses the median NN, so there is a further $\sim$30\% (larger in the more clustered regions) slippage between the reported median and the strict line-density scale; we carry this as an additional stated systematic, of the same order as the convention systematics, rather than fold it in --- it does not affect the sub-classical \emph{sign} (it makes the mean-gap line density scale even shorter) but it should be remembered in any order-of-magnitude comparison with a simulated wavelength. The observed median separation is therefore essentially set by the core line-density plus a finite-core exclusion, \emph{not} by a resonant fragmentation wavelength: the ``spacing'' is sub-classical in the sense that cores are more numerous along the spine than one per classical wavelength ($\lambda_{\rm m}\approx5$--$6\times$FWHM would place $\sim3$--$4\times$ fewer cores), the signature of sub-fragmentation rather than a single characteristic scale.
\begin{figure}
\centering
\includegraphics[width=0.99\columnwidth]{figures/fig_clustering.pdf}
\caption{Adjacent-core gap distributions along the skeletons of the four robust regions (histograms), with exponential/Poisson (solid) and regular/periodic (dashed) reference curves; each panel is annotated with the coefficient of variation CoV and the 1-D Clark--Evans ratio $R$. CoV of order unity and the overwhelming rejection of the periodic model ($p\lesssim10^{-15}$) show the cores are random-to-clustered, not periodically spaced.}
\label{fig:clustering}
\end{figure}

This does not make the comparison with a fragmentation wavelength meaningless, but it fixes its logical status. A fragmentation process still sets a characteristic \emph{mean} core separation --- the number of cores formed per unit length --- even when the resulting point pattern is not a sharp periodic comb; clustering, a finite-core (beam-scale) exclusion, and projection across blended structure then broaden that mean into the observed near-exponential gap distribution. The observed median separation and any single simulated fragmentation wavelength therefore need not, and should not, coincide better than order-of-magnitude: what is physical and convention-robust is that the core \emph{line density} is sub-classical --- roughly three to four times more cores per unit length than the one-core-per-$\lambda_{\rm m}$ of the classical equilibrium prediction. We accordingly read all wavelength-vs-spacing comparisons below as order-of-magnitude consistency of a characteristic scale, not as a match of a resonant wavelength.

\subsection{Combined Systematic Error Budget}
\label{sec:error_budget}

We distinguish three kinds of uncertainty and present them separately rather than as a single error bar: (i) \emph{statistical} (bootstrap/jackknife resampling, formal); (ii) \emph{systematic} (measured offsets, e.g.\ the pipeline bias and the fixed-width convention); and (iii) \emph{sensitivity ranges} (plausible bounds on effects we can sign but not measure as Gaussian errors, e.g.\ protostellar migration). Quoted $\pm$ values combine these in quadrature for economy but are \emph{not} formal statistical confidence intervals (see also the Abstract). We further separate the observational error budget of each statistic from the additional systematics that enter only when simulation predictions are compared with observation. For the nearest-neighbour primary measurement, the primary region-Plummer NN result $\lambda/W\approx1.4$ (raw; the raw fixed-width analogue is $\approx2.1$) carries four observational systematics: residual pipeline bias ($\sim$10\%, from the spread in injection--recovery ratios 1.11--1.29 after bias correction; the uncorrected pipeline bias is $\sim$15--20\%, and the intrinsic NN statistic bias is 3\%); the fixed-width convention ($\sim$25\%; the distance-dependent filament-width variation measured by \citealt{Panopoulou2022}, including its beam-convolution component, counted once here, not in addition to it); protostellar migration ($\sim$7.5\%, uncorrected); and between-region bootstrap scatter ($\sim$4\%). Manual spur removal during skeleton cleaning is a further, unquantified systematic, smaller than the width-convention term but not bounded by independent re-editing. The skeleton-extraction algorithm is the single largest observational systematic, an independent FilFinder pipeline giving values $2$--$3\times$ smaller than DisPerSE (Section~\ref{sec:injection_recovery}); we treat it separately below. The dominant term (the width convention) is a \emph{deterministic} offset, not a random error, so we regard the aggregate not as a Gaussian $\sigma$ but as a \emph{min--max envelope}: the primary value spans $\lambda/W\approx0.8$--$1.9$ across the plausible convention and correction choices (Table~\ref{tab:headline}), which is the number a reader should carry. For economy we also quote a single quadrature figure, $\sigma_{\rm sys}\approx28\%$, i.e.\ a $\sim28\%$ systematic on the primary value (equivalently $\lambda/W=1.9\pm0.5$ under the fixed-width convention), with the skeleton factor of two carried separately.

The legacy pairwise-median value $\lambda/W=2.84$ is an observational $L/3$ scale measure (Section~\ref{sec:statistics}); its observational uncertainty is the bootstrap scatter $\pm4.2\%$ ($\lambda/W=2.84\pm0.12$), and we do not quote a $\sigma$-level tension with the classical prediction because the statistic does not measure the fragmentation wavelength. (In pc, $\lambda/W=2.84$ is $0.284$\,pc; the full- and robust-sample core-weighted pairwise means are $0.279$ and $0.290$\,pc respectively, Table~\ref{tab:sample}, differing from $0.284$ only through the weighting scheme.)

\emph{Skeleton-algorithm dependence.} The factor-of-two spread between DisPerSE and FilFinder is the single largest uncertainty, and the two algorithms do not measure the same object. DisPerSE identifies topologically robust density crests via persistence, tracing the ridge of each filament, and is the algorithm used to build the published HGBS filament catalogues, so it is the natural choice for consistency with prior work. FilFinder instead skeletonises an intensity mask, returning the medial axis of a thresholded region; at HGBS resolution this yields a denser, more branched network that subdivides a single filament into several segments, placing more spine near each core and so shortening the apparent nearest-neighbour spacing. The two bracket a genuine ambiguity in what constitutes ``the filament'' at these scales, rather than one being simply wrong, and we regard their values ($\lambda/W\approx1.9$ and $\approx0.4$--$0.8$ respectively) as systematic limits on the true spacing. The point is convention-independent: both lie far below the classical $\lambda_{\rm m}\approx5$--$6\times$FWHM, so the sub-classical result is not an artefact of the skeletoniser, even if its precise value is. A definitive resolution, common-footing extraction with matched spatial filtering and an agreed definition of a filament segment, is beyond our scope, and we flag the factor of two as an open observational systematic rather than resolve it here. Until a community consensus exists on filament definition and skeleton extraction, the observational core spacing cannot be known to substantially better than a factor of $\sim$2, and this, rather than the formal statistical error, is the dominant limit on any comparison with theory. The choice of skeletoniser also selects the physical interpretation. In matched (beam-convolved Plummer) width units the simulated single-filament wavelength is $\approx1.0$--$1.3$ (Section~\ref{sec:validation_endtoend}). It coincides with the region-Plummer observed values ($1.2$--$1.4$, raw), sits \emph{below} the fixed-width DisPerSE spacing ($1.9$, periodic-corrected; single filaments slightly over-produce fragmentation relative to it), and lies \emph{above} the FilFinder spacing ($0.4$--$0.8$). So if the denser FilFinder extraction is the more faithful one, single-filament fragmentation would be insufficient and the extra compression of hierarchical-fibre projection (Section~\ref{sec:fibre_projection}) would be required; under DisPerSE or the region-Plummer widths it is not. The two skeletonisers thus bracket not only the spacing but which explanation it favours, and only a fibre-resolved measurement (Section~\ref{sec:limitations}) can decide between them.

Table~\ref{tab:errorbudget} collects every significant uncertainty on the comparison. The dominant terms are \emph{methodological}, not statistical: the skeleton algorithm and the width convention each move $\lambda/W$ by up to a factor of two, and it is these, rather than the formal bootstrap error, that limit how precisely the observed and simulated spacings can be compared.

\begin{table}
\caption{Summary of the uncertainties on the $\lambda/W$ comparison. ``Sim.'' entries apply to the simulated wavelength, ``obs.'' to the observed spacing, and ``comparison'' to terms that enter only when the two are matched.}
\label{tab:errorbudget}
\footnotesize
\setlength{\tabcolsep}{3pt}
\begin{tabular}{p{3.6cm}p{1.9cm}p{1.7cm}}
\toprule
Source & Type & Magnitude \\
\midrule
Bootstrap / cloud-to-cloud (obs.) & sensitivity/descriptive & $\sigma\approx0.4$ ($\sim$20\%) \\
NN pipeline bias (obs.) & systematic & $\sim$10\% (corrected) \\
Protostellar migration (obs.) & sensitivity & $\sim$7.5\% \\
Width convention, fixed vs region (obs.) & systematic & $\sim$25\% \\
Skeleton algorithm, DisPerSE/FilFinder (obs.) & systematic & factor $\sim$2 \\
Growth-rate resolution/epoch (sim.) & modelling & $\lesssim$5\% (converged) \\
Line-mass ($f$) dependence (sim.) & modelling & $\sim$25\% \\
Gaussian vs Plummer FWHM (comparison) & systematic & factor $\sim$1.5 \\
\bottomrule
\end{tabular}
\end{table}

When simulation $\lambda/W_{\rm core}$ values are converted to observable $\lambda/W_{\rm fil}$, the T1 width-normalisation (forward-modelled $\eta_W=0.59$--$0.78$, provisional calibration $0.61\pm0.10$; Section~\ref{sec:width_correction}) is the dominant additional systematic; the directly measured supercritical ensemble (Section~\ref{sec:modecompetition}) carries its own $\pm0.4$ scatter. These comparison-only systematics do not affect the NN observational measurement, which is already in $W_{\rm fil}$ units.

Because the four robust clouds are the genuine independent units, we quote uncertainties descriptively, not inferentially. Two distinct quantities must be kept apart: (i) the empirical \emph{per-region range} of the bias-corrected fixed-width values, $\lambda/W=1.0$--$2.1$ (Taurus, Orion~B, Perseus, Aquila $=1.0,1.9,2.1,1.8$; Section~\ref{sec:injection_recovery}), which is the real cloud-to-cloud spread; and (ii) the \emph{methodological/systematic envelope}, $\lambda/W\approx1.1$--$2.4$, obtained by varying the width-convention, skeleton and association-radius choices (the quadrature sum of these is $\sigma_{\rm sys}\approx28\%$, a synthesis of heterogeneous sensitivity ranges, not a Gaussian error). With $N=4$ we do \emph{not} report a bootstrap $\sigma$ or a confidence interval as a formal quantity, and we do not claim a $\sigma$-level exclusion of the classical prediction; the informal descriptor ``$\sigma\approx0.4$'' quoted elsewhere for the four-cloud scatter should be read only as the half-width of the per-region range.

\subsection{NN pipeline reproducibility, population splits and hierarchical statistics}
\label{sec:nnpipeline}

The along-skeleton NN pipeline thresholds the deposited DisPerSE skeleton, bridges fragmentation by morphological closing and re-skeletonisation, associates each core to the nearest spine within $0.06$\,pc, orders cores along each component by graph arclength (double-BFS diameter), and takes adjacent-core separations. It reproduces the deposited catalogue counts and raw medians exactly (Table~\ref{tab:nnpipeline}); cores not on a robust skeleton segment are rejected.

\begin{table}
\caption{NN pipeline counts: skeleton components, catalogue cores, cores associated to a skeleton, and adjacent NN pairs, per robust region.}
\label{tab:nnpipeline}
\footnotesize
\setlength{\tabcolsep}{4pt}
\begin{tabular}{lcccc}
\toprule
Region & comps & catalogue & associated & NN pairs \\
\midrule
Orion B & 67 & 1844 & 680 & 614 \\
Aquila  & 14 & 749  & 182 & 168 \\
Perseus & 54 & 816  & 619 & 563 \\
Taurus  & 32 & 536  & 479 & 461 \\
\bottomrule
\end{tabular}
\end{table}

The fraction of catalogue cores that associate to a robust skeleton segment varies substantially across the four regions --- Orion~B $37\%$ ($680/1844$), Aquila $24\%$ ($182/749$), Perseus $76\%$, Taurus $89\%$ --- so in Orion~B and Aquila the NN spacing is measured on a minority of the cores. This spread reflects skeleton \emph{coverage}: DisPerSE deposits crests only where the persistence-filtered column density exceeds threshold, so the more distributed, clustered clouds (Orion~B, Aquila) have many cores in inter-filament or low-contrast material that lies off any robust crest, whereas the filament-dominated clouds (Perseus, Taurus) place almost all cores on a spine. Any bias from this, if present, runs \emph{toward} the sub-classical result, not against it: cores rejected because they lie on real but sub-threshold (fainter, lower-line-density) segments carry \emph{longer} spacings, so omitting them biases the measured median \emph{short}. We therefore bound its size directly rather than merely signing it. The decisive check is that the result does \emph{not} depend on the low-coverage regions: restricting to the two \emph{high}-coverage clouds (Perseus $76\%$ and Taurus $89\%$, which capture essentially all on-filament cores) gives $\lambda/W=1.8$ and $1.1$, median $1.4$ --- identical to the four-region median --- so the low $24$--$37\%$ coverage in Orion~B/Aquila cannot be driving the sub-classical characteristic. Two further checks agree: the Orion~B/Aquila NN medians are stable to $\lesssim5\%$ as the association half-width varies over $0.04$--$0.08$\,pc, and the prestellar subset (which dominates the on-skeleton budget) gives the same headline as the all-cores measurement (below). A dedicated injection--recovery test dropping the faintest skeleton segments would sharpen the bound further; the region-restricted bound already shows the effect is not load-bearing.

To test whether mixing core populations biases the headline, we repeated the measurement separately for prestellar, unbound starless, and protostellar cores using the published classification flags. The prestellar-only robust-median $\lambda/W=2.10$ ($717$ pairs) is statistically indistinguishable from the all-cores value $2.08$ ($1806$ pairs; both $\approx1.9$ after the injection-recovery bias correction). Prestellar cores dominate the on-skeleton budget in every region and set the headline; unbound starless cores give a larger, noisier value (median $\lambda/W\approx3.3$) and protostellar cores are too few ($6$--$80$ pairs) to constrain it. The sub-classical headline is therefore not an artefact of population mixing.

\smallskip\noindent\textbf{Hierarchical statistics and angular-resolution robustness.}

Cores within a filament are not independent, so we assessed the role of the cloud-to-filament-to-core nesting with a three-level nonparametric bootstrap (clouds $\to$ skeleton-connected filament components $\to$ core spacings, 2000 iterations; adjacent spacings sharing a central core are carried in and out together with their filament, capturing the dominant covariance). On a 2-D projected nearest-neighbour metric, which underestimates the along-spine arc length and so sits below the headline value ($\lambda/W\approx1.1$--$1.5$), the bootstrap median is stable and does not shift the central value; the dominant source of uncertainty is the intrinsic cloud-to-cloud scatter, $\sigma\approx0.4$ in $\lambda/W$ (range $\sim0.8$ in Taurus to $\sim1.7$ in Aquila). Core-count and equal-cloud weighting agree to within $\sim0.1$. We therefore quote the cloud-to-cloud scatter, rather than a per-pair formal error, as the relevant uncertainty on the characteristic spacing (Section~\ref{sec:error_budget}).

Angular-resolution blending: the HGBS beam grows from $0.012$\,pc (Taurus, 135\,pc) to $0.039$\,pc (Aquila, 436\,pc), and the fraction of nearest-neighbour pairs closer than $2\times$ the physical beam rises with distance ($0$\,per cent in Taurus to $10$\,per cent in Orion~B; Spearman $\rho=+0.80$), as expected. Merging cores within one beam shifts the characteristic $\lambda/W$ by $\le0.004$, however, so distance-dependent blending is not a significant bias on the headline.

\begin{table}
\caption{Summary of validation tests performed in this paper.}
\label{tab:validation_summary}
\footnotesize
\setlength{\tabcolsep}{3pt}
\begin{tabular}{p{2.4cm}p{2.8cm}p{2.0cm}}
\toprule
Test & What it checks & Result \\
\midrule
Injection--recovery (NN) & NN bias on real skeletons & $3\%$; pipeline $15$--$20\%$ \\
Injection--recovery (pairwise) & $L/3$ convergence & tracks $L/3$, not $\lambda$ \\
Resolution convergence & $t_{\rm frag}(128^3)$ vs $256^3$ & ratio $0.915\pm0.012$ \\
Truelove criterion & cells per $\lambda_{\rm J}$ at measurement & long.\ $5$--$7$; perp.\ $1$--$2$ \\
IC sensitivity & Gaussian vs uniform profile & identical \\
EOS sensitivity & $\gamma=0.8$--$5/3$ & $\gamma<1$ accelerates; $5/3$ suppresses \\
Region resampling (descriptive) & spread over 4 clouds & $0.261$--$0.298$\,pc \\
Leave-one-out & single-region dominance & $<6\%$ shift \\
Population splits & prestellar vs all cores & $\lambda/W\!=\!2.10$ vs $2.08$ \\
Hierarchical bootstrap & cloud$\to$filament$\to$core & $\sigma_{\rm cloud}\!\approx\!0.4$ \\
Mock-blending & distance-dependent beam & $\Delta\lambda/W\!\le\!0.004$ \\
Region-specific Plummer widths & beam-deconvolved $W_{\rm fil}$ & $0.11$--$0.16$\,pc; $\lambda/W\!=\!1.1$--$1.8$ \\
\bottomrule
\end{tabular}
\end{table}

\subsection{Comparison with Literature}

Our revised-distance measurements differ from the original HGBS values (Table~\ref{tab:literature}) by the region-dependent distance revisions (Aquila $+68\%$; Orion~B essentially unchanged, its pre-Gaia distance already $\sim400$\,pc; Taurus $-4\%$); the core-count-weighted mean rises only $\sim10\%$ for the robust regions, dominated by Aquila.

\begin{table*}
\caption{Comparison with Published HGBS Spacing Measurements}
\label{tab:literature}
\begin{tabular}{lccccc}
\toprule
Region & This Work, NN (pc) & Original HGBS (pc) & This Work, pairwise (pc) & Literature (pc) & Reference \\
\midrule
\textbf{Robust Regions} &  &  &  &  &  \\
Orion B & \textbf{0.207} & 0.212 & 0.313 $\pm$ 0.047 & 0.20--0.30 & \citet{Konyves2020} \\
Aquila  & \textbf{0.210} & 0.206 & 0.346 $\pm$ 0.047 & 0.22--0.26 & \citet{Konyves2015} \\
Perseus & \textbf{0.263} & 0.199 & 0.248 $\pm$ 0.040 & 0.22 $\pm$ 0.02 & \citet{Pezzuto2021} \\
Taurus  & \textbf{0.124} & 0.206 & 0.198 $\pm$ 0.040 & 0.19 $\pm$ 0.02 & \citet{Hacar2013} \\
\midrule
\textbf{Limited Regions} &  &  &  &  &  \\
Ophiuchus & --- & 0.197 & 0.206 $\pm$ 0.053 & 0.18 $\pm$ 0.02 & \citet{Konyves2020} \\
Serpens  & --- & 0.188 & 0.331 $\pm$ 0.097 & 0.17--0.19 & \citet{Konyves2020} \\
TMC1     & --- & 0.203 & 0.195 $\pm$ 0.056 & $\sim$0.20 & \citet{Hacar2013} \\
CRA      & --- & 0.212 & 0.248 $\pm$ 0.072 & $\sim$0.21 & \citet{Konyves2020} \\
\bottomrule
\end{tabular}

\vspace{0.1in}
\footnotesize
\emph{Notes} (1) The primary statistic is the \emph{nearest-neighbour} (NN) spacing (bold). On this statistic the homogeneous re-analysis leaves Orion~B ($0.207$ vs $0.212$\,pc) and Aquila ($0.210$ vs $0.206$\,pc) essentially unchanged, \emph{shortens} Taurus ($0.124$ vs $0.206$\,pc; highest association fraction, so most reliable), and \emph{lengthens} Perseus ($0.263$ vs $0.199$\,pc); it does \emph{not} inflate the spacing the way the pairwise median does, and the net effect strengthens rather than revises away the sub-classical result. The ``This Work, pairwise'' column is the all-pairs pairwise median, retained only as a like-for-like comparator with legacy pairwise measurements; it is inflated towards the region extent $L/3$ (Section~\ref{sec:statistics}) and should not be read as a physical spacing. NN spacings for the limited regions are not used and are omitted. Perseus is the one region where the along-skeleton NN ($0.263$) slightly \emph{exceeds} the pairwise median ($0.248$), the reverse of the other three: on Perseus's more fragmented, heavily bridged skeleton the arclength ordering (which follows the curved, reconstructed spine) inflates adjacent-core \emph{path} distances relative to straight-line pairs, so the two statistics cross; both remain sub-classical, and Perseus is anyway the least sub-classical robust region. (2) ``Original HGBS'' values are the published spacing measurements using pre-Gaia distances. (3) Literature ranges reflect different measurement methods (pairwise vs.\ nearest-neighbour) and distance assumptions. (4) Serpens original HGBS value used $260$\,pc; revised $458$\,pc ($+76\%$) increases the pairwise spacing proportionally. (5) For Orion~B the pre-Gaia HGBS distance was already $\sim$400\,pc \citep{Konyves2020}, so the pairwise $0.212\to0.313$\,pc difference is a \emph{methodological} (statistic-definition) effect of the homogeneous re-analysis, not a distance revision.
\end{table*}

We decomposed this shift on the real Orion~B data by toggling each methodological choice in turn at fixed distance. The dominant term by far is the \emph{statistic definition}: switching from the nearest-neighbour separation to the per-filament all-pairs pairwise median inflates the spacing by a factor $\approx1.5$ ($0.207\to0.313$\,pc), because the pairwise median is biased towards the region extent ($L/3$; Section~\ref{sec:statistics}) rather than measuring a spacing, and this single change accounts for essentially the entire $0.212\to0.313$\,pc difference in Table~\ref{tab:literature}. Skeleton re-extraction/bridging ($+4\%$), association half-width ($\pm12\%$) and 2-D$\to$arclength ordering ($+12\%$) each contribute at the $\le15\%$ level and largely cancel, and the distance revision ($400\to386$\,pc) contributes $-3.5\%$. Crucially, on our \emph{primary} nearest-neighbour statistic Orion~B is $0.207$\,pc --- within $3\%$ of the \citet{Konyves2020} value --- so the apparent shift is a property of the pairwise-median comparability statistic, not of the physical spacing. This like-for-like NN comparison is not uniform across the four robust regions, and we state the full picture rather than the best case: the homogeneous NN re-analysis leaves Orion~B ($0.207$ vs $0.212$\,pc) and Aquila ($0.210$ vs $0.206$\,pc) essentially \emph{unchanged}, \emph{shortens} Taurus ($0.124$ vs $0.206$\,pc; not a distance effect, and measured at the highest association fraction, $89\%$, so arguably the most reliable region), and \emph{lengthens} Perseus ($0.263$ vs $0.199$\,pc). The net effect on the sub-classical conclusion is if anything to strengthen it (Taurus, the most reliable region, moves shorter); no region's NN spacing is inflated in the way the pairwise median is.

\section{THEORETICAL FRAMEWORK}
\label{sec:theory}

\subsection{Classical Fragmentation Theory}

For an isothermal, self-gravitating cylinder of density $\rho_0$ and sound speed $c_s$, the critical line mass for gravitational instability is \citep{Ostriker1964}:
\begin{equation}
\mu_{\rm crit} = \frac{2c_s^2}{G}
\end{equation}
The most unstable fragmentation wavelength is \citep{Inutsuka1992}:
\begin{equation}
\lambda_{\rm IM92} = 22\,H \approx 4\times(2R_{\rm flat}) \approx 5\text{--}6\times W_{\rm fil}
\label{eq:im92}
\end{equation}
where $H = c_s\,(4\pi G\rho_c)^{-1/2}$ is the thermal scale height ($\rho_c$ the central density), $R_{\rm flat}=\sqrt8\,H$ is the flattening radius, and $W_{\rm fil}\approx0.1$\,pc is the characteristic filament width, taken throughout as the FWHM of the (beam-deconvolved) column-density profile. The classical prediction is most reliably quoted in FWHM units: solving the dispersion relation exactly (Appendix~\ref{sec:val_appendix}; \citealt{Fischera2012}) gives a fastest-growing wavelength $\lambda_{\rm m}\approx5\times W_{\rm fil}$ for the near-critical, star-forming regime, rising to $\approx6\times W_{\rm fil}$ at low line mass ($f_{\rm cyl}\to0$; \citealt{Fischera2012}), and a shortest-unstable (critical) wavelength $\lambda_{\rm cr}\approx2.6\times W_{\rm fil}$. The often-quoted ``$\lambda\approx4\times$ the width'' is $4\times$ the flat \emph{diameter} $2R_{\rm flat}$ (\citealt{Inutsuka1992}), which equals $\approx5$--$6\times$ the \emph{FWHM}; taking it as $4\times$ the FWHM conflates the two and \emph{understates} the classical scale. We therefore adopt $\lambda_{\rm m}\approx5$--$6\,W_{\rm fil}$ (Fischera \& Martin 2012) as the benchmark. This ratio is intrinsically defined on the FWHM of the classical \emph{equilibrium} ($p=4$ Ostriker) column-density profile, whereas the observed filaments are better fit by $p\simeq2$; because we compare in FWHM units (a directly measured width for either profile) and, for the simulations, forward-model to a beam-convolved Plummer FWHM measured exactly as on the data (Section~\ref{sec:width_correction}), the sub-classical \emph{deficit} does not depend on the profile-index difference --- the observed cores are closer than the classical instability predicts however the width is parametrised. Because observed and simulated spacings and widths are measured at the same epoch, we compare them as the density-independent ratio $\lambda/W_{\rm fil}$, which sidesteps the fact that $22H\propto\rho_c^{-1/2}$ is not fixed in absolute units. The simulation width bookkeeping is set out in Section~\ref{sec:width_correction}. This sub-classical result is not new in sign: \citet{Fischera2012} already noted that observed cores sit ``about the FWHM'' apart, too close for the classical instability, and \citet{Zhang2023cal} find a core spacing $\approx1.3\times$ the width in the California cloud, significantly below the canonical value. Our contribution is to quantify the deficit homogeneously with Gaia DR3 distances and to reproduce it in a controlled MHD campaign.

With $f=\mu_{\rm line}/\mu_{\rm crit}$ and $t_{\rm frag}\propto(f-1)^{-1/2}$, marginally supercritical filaments fragment slowly. At magnetically supercritical field strength ($\beta\gtrsim0.2$) all configurations fragment for $\mathcal{M}\geq1$ at every $f$ tested once integrated adequately: above the critical line mass there is no \emph{thermal} stability boundary. This is distinct from \emph{magnetic} stabilisation, where a strong longitudinal field ($\beta\lesssim0.15$) suppresses collapse even for $f>1$ (the ``strong-field, magnetically supported'' regime; Appendix~\ref{sec:supplementary_campaigns_app}); the thermal ($f$) and magnetic (mass-to-flux) criticalities are physically distinct.

We adopt a single reference $W_{\rm fil}=0.10$\,pc for all regions, for comparability with prior HGBS work. The universality of this width is debated \citep{Panopoulou2017,Panopoulou2022}. We have measured region-specific Plummer widths directly (Section~\ref{sec:widths}); they are $0.11$--$0.16$\,pc ($p\simeq2.0$) and give $\lambda/W\approx1.4$, so the sub-classical conclusion is robust to the width convention, and the width term is carried once as the single $\sim25\%$ systematic in the Section~\ref{sec:error_budget} budget.

\subsection{Width normalisation: $W_{\rm core}$ to $W_{\rm fil}$}
\label{sec:width_correction}
Three widths enter the simulation--observation comparison and must be kept separate (Table~\ref{tab:widthdefs}): the initial Gaussian half-width $W_{\rm core}=0.3\,\lambda_{\rm J}$; the formation FWHM $W_{\rm form}=0.683\,\lambda_{\rm J}$ of the evolved projected profile ($\simeq2.28\,W_{\rm core}$, slightly below the initial $2.355\,W_{\rm core}$ owing to contraction; resolution- and seed-converged to $<1\%$); and the observed HGBS width $W_{\rm fil}$, the beam-convolved Plummer FWHM \citep{Arzoumanian2019}, which we obtain by forward-modelling the simulated filaments through the full synthetic-HGBS pipeline (projection, beam convolution, Plummer fit). This gives $\eta_W=W_{\rm form}/W_{\rm fil}=0.59$--$0.78$ and
\begin{equation}
T_1 \equiv \frac{W_{\rm form}}{W_{\rm fil}} = 0.61 \pm 0.10 \quad (\text{forward-model band } 0.59\text{--}0.78),
\label{eq:t1_correction}
\end{equation}
so that $W_{\rm fil}=W_{\rm form}/T_1\approx1.12\,\lambda_{\rm J}$. The observable is $\lambda/W_{\rm fil}$; because the fragmentation spacing $\lambda$ is a physical length independent of the width definition, the conversion from the simulation-internal ratio $\lambda/W_{\rm core}$ (normalised by the initial $\sigma=0.3\,\lambda_{\rm J}$) is
\begin{equation}
\left(\frac{\lambda}{W}\right)_{\rm fil} = \frac{W_{\rm core}}{W_{\rm fil}}\left(\frac{\lambda}{W}\right)_{\rm core} = T_1\,\frac{W_{\rm core}}{W_{\rm form}}\left(\frac{\lambda}{W}\right)_{\rm core} \approx 0.27\left(\frac{\lambda}{W}\right)_{\rm core},
\label{eq:t1_apply}
\end{equation}
i.e.\ $\lambda/W_{\rm fil}=\lambda\,[\lambda_{\rm J}]/1.12$, the spacing in units of the forward-modelled observed width; the factor $0.27$ includes the $\sigma\!\to\!$FWHM factor $0.44$, and using $T_1$ alone would overstate $\lambda/W_{\rm fil}$ by $\approx2.3\times$. The observational NN measurement (Section~\ref{sec:injection_recovery}) is already in $W_{\rm fil}$ units. Because the observable is sensitive to how $W_{\rm fil}$ is defined, we do not rely on it; the primary comparison uses the convention-free growth-rate wavelength (an upper limit; Section~\ref{sec:validation_endtoend}), and this normalisation is retained only as a cross-check.

\begin{table}
\caption{The three filament widths used in the simulation--observation comparison (all in Jeans lengths, $\lambda_{\rm J}=1$).}
\label{tab:widthdefs}
\footnotesize
\setlength{\tabcolsep}{3pt}
\begin{tabular}{p{1.5cm}p{1.4cm}p{4.0cm}}
\toprule
Symbol & Value ($\lambda_{\rm J}$) & Definition \\
\midrule
$W_{\rm core}$ & $0.30$ & initial Gaussian half-width $\sigma$ (input) \\
$W_{\rm form}$ & $0.683$ & Gaussian FWHM of evolved projected profile, supercritical ensemble ($\simeq2.28\,W_{\rm core}$) \\
$W_{\rm FWHM}$ & $0.58$ & Gaussian FWHM at the growth-rate epoch, near-critical runs (narrower, already contracting); normalises Table~\ref{tab:wave} \\
$W_{\rm fil}$ & $\approx1.12$ & beam-convolved Plummer FWHM ($=W_{\rm form}/T_1$); the observable normalisation \\
\bottomrule
\end{tabular}
\end{table}

For the ambient-confined supercritical ensemble ($\lambda/W_{\rm core}=3.27$) this gives $\lambda/W_{\rm fil}\approx0.9$.

\subsection{Magnetic Tension Mechanism}
\label{sec:magnetic}

For a filament threaded by a magnetic field at an angle $\theta$ to the axis (Alfv\'en speed $v_A$), the general oblique dispersion relation is \citep{Nakamura1993}:
\begin{equation}
\omega^2 = c_s^2 k^2 + v_A^2 k^2 \sin^2\theta - 4\pi G\rho_0 F(kR)
\label{eq:dispersion}
\end{equation}
where $F(kR)$ encodes the cylindrical geometry and $\theta$ is the field--axis angle ($0^\circ$ longitudinal, $90^\circ$ perpendicular). This is the thin-perturbation \emph{static-field} dispersion relation: it treats $\mathbf{B}$ as a fixed background threading the cylinder and by construction does \emph{not} describe a current-carrying (helical) equilibrium, which is precisely why the helical/toroidal case cannot be read off from Eq.~(\ref{eq:dispersion}) and requires the separate treatment of Section~\ref{sec:optional_tests}. For the purely axial mode of a \emph{longitudinal} field ($\theta=0^\circ$) the compression is directed along $\mathbf{B}$, which exerts no restoring force, so the magnetic term $v_A^2k^2\sin^2\theta$ vanishes identically and Eq.~(\ref{eq:dispersion}) reduces to the field-free thermal relation: a static longitudinal field leaves the fragmentation wavelength almost unchanged. (The magnetic term is maximal for a perpendicular field, $\theta=90^\circ$, and it is this term that would enter a perpendicular magnetic-Jeans calibration.) We do \emph{not} solve the full cylindrical dispersion relation, and this paper therefore contains \emph{no quantitative model of magnetic tension}. Consistent with the physics above, our longitudinal-field runs show the \emph{growth-rate} wavelength depends only weakly on $\beta$ ($\lambda_{\rm max}=1.00,1.14,0.80\,\lambda_{\rm J}$ at $\beta=0.3,0.5,1.0$, a non-monotonic spread within the axial mode-quantisation step $\Delta\lambda\approx0.15\,\lambda_{\rm J}$ set by $L_x=8\,\lambda_{\rm J}$, so not a resolved physical trend; the larger $\beta$-trend of the late-time bead count (a resolution-sensitive, late-epoch measure we do not adopt for the wavelength; Appendix~\ref{sec:supplementary_campaigns}) is an artefact of that epoch) and on supercriticality, so \emph{in our runs} the scale is set by the thermal Jeans length, not the field. The shortening we measure is therefore \emph{not} a magnetic-tension effect: it is the dynamical/contraction shift. A genuinely shorter magnetic wavelength would require a \emph{helical}/toroidal field in force balance; our seeded helical runs (Section~\ref{sec:optional_tests}) narrow the width by a radial pinch but do not shorten $\lambda_{\rm max}$, and the force-balanced helical equilibrium in which the \citet{Fiege2000} pinch mode operates is left to future work. Planck polarimetry finds dense filaments preferentially perpendicular to the mean \emph{projected} field \citep{Jadhav2026}; the projected orientation does not fix the internal three-dimensional geometry, so the internally-longitudinal fraction is strictly unconstrained --- if it tracks the projected orientation it is a minority ($\lesssim10\%$), which is the assumption under which our longitudinal result applies to only part of the population. These estimates are calibrated in the near-critical regime ($f\approx1$); the directly measured supercritical spacing appears in Section~\ref{sec:modecompetition}.

\subsection{Hierarchical Fragmentation}

High-resolution studies have revealed that filament fragmentation is hierarchical: filaments fragment into velocity-coherent fibres, and fibres fragment into cores \citep{Hacar2013,Hacar2018,Yang2024}. Whether the fibre-to-core spacing is quasi-periodic at the classical scale is not yet established, the ALMA study of \citet{Yang2024} finds a clear filament--fibre--core hierarchy but only a weak quasi-periodic imprint at the core level, so the fibre-level fragmentation scale remains an observational open question.

If HGBS filaments are bundles of velocity-coherent fibres, each fragmenting at the classical scale, then measuring core spacings across the whole filament rather than within individual fibres would compress the apparent spacing. This goes in the observed direction: fibre-level measurements recover the classical $\lambda_{\rm m}\approx5$--$6\times$FWHM prediction, while filament-level measurements are sub-classical.

We test this interpretation quantitatively in Section~\ref{sec:fibre_projection}, where a forward-model of $N$-fibre bundles run through our identical NN pipeline reproduces the observed $\lambda/W$; it emerges there as a viable explanation for the sub-classical filament-level spacing (though, as Section~\ref{sec:validation_endtoend} shows, not the only one, since single-filament dynamical fragmentation is also consistent with the data). The evidence is not yet decisive: the \citet{Yang2024} fibre-level result comes from a single region (Orion~B) and may not generalise, and the number of fibres per filament is not directly constrained across our robust regions. Confirming the picture requires fibre-resolved core-spacing analysis in further HGBS filaments, to test whether the classical $\lambda_{\rm m}\approx5$--$6\times$FWHM recovers within individual fibres beyond Orion~B.

\section{MHD SIMULATIONS}

The full campaign comprises $2251$ Athena++ runs across 18 grids (Table~\ref{tab:campaigns}); it is a set of controlled numerical experiments, not models of individual HGBS filaments. Most of these runs establish \emph{outcomes and timescales} across parameter space; the headline fragmentation-\emph{wavelength} measurement rests on the 25-run near-critical growth-rate campaign (whose dominant-mode measurements across $f$, $\beta$ and resolution are reported in Table~\ref{tab:wave}; the 8 isotropic $128^3$--$512^3$ runs at $f=1.1$ form the production convergence ladder), so the large campaign size should not be read as statistical weight behind the wavelength. Two results emerge (Section~\ref{sec:why_mhd_fails}). The \emph{outcome} of a run, radial collapse or longitudinal beading, depends on regime and boundary conditions, so a single filament has no setup-independent fate. Its dominant fragmentation \emph{mode}, measured from the growth-rate spectrum, is sub-classical ($\approx0.5$--$1\times$ the classical value; Section~\ref{sec:validation_endtoend}) for near-critical, longitudinal-field filaments; the perpendicular geometry is not numerically resolved. This section provides the evidence for both.

\subsection{Simulation Methodology}

We performed self-gravitating MHD simulations using Athena++ \citep{Stone2020} to test filament fragmentation under controlled conditions. The simulations solve the ideal MHD equations with self-gravity:
\begin{align}
\frac{\partial \rho}{\partial t} &= -\nabla \cdot (\rho \mathbf{v}) \\
\frac{\partial (\rho \mathbf{v})}{\partial t} &= -\nabla \cdot \left(\rho \mathbf{v} \mathbf{v} + P + \frac{B^2}{8\pi} - \frac{\mathbf{B}\mathbf{B}}{4\pi}\right) - \rho \nabla \phi \\
\frac{\partial \mathbf{B}}{\partial t} &= \nabla \times (\mathbf{v} \times \mathbf{B}) \\
\frac{\partial E}{\partial t} &= -\nabla \cdot \left[(E + P + \frac{B^2}{8\pi})\mathbf{v} - \frac{(\mathbf{v} \cdot \mathbf{B})\mathbf{B}}{4\pi}\right] - \rho \mathbf{v} \cdot \nabla \phi
\end{align}
where $\rho$ is density, $\mathbf{v}$ is velocity, $\mathbf{B}$ is magnetic field, $P$ is pressure, $E$ is total energy, and $\phi$ is gravitational potential satisfying $\nabla^2 \phi = 4\pi G\rho$.

We initialise a cylindrical filament with density profile $\rho(r) = \rho_0 / [1 + (r/W)^2]$ and uniform magnetic field $\mathbf{B}_0$. The filament is characterised by: the thermal line-mass fraction $f = M_{\rm line}/M_{\rm line,crit}$ (Eq.~2), plasma $\beta = 8\pi P/B^2$ (for reference, at filament conditions $n\sim10^4\,{\rm cm^{-3}}$ and $T\sim10$\,K, $\beta=1$ corresponds to $B\sim15\,\mu$G and $\beta=0.3$ to $\sim30\,\mu$G, i.e.\ the $\mu$G--tens-of-$\mu$G range of Zeeman and dust-polarimetry estimates in dense filaments; \citealt{Crutcher2012}), and the seed-amplitude parameter $M_{\rm seed}=\delta v/c_s\times10^{4}$ (a perturbation-spectrum label, not a physical Mach number; the base outcome/timescale grid uses a single-mode seed $\delta v/c_s=10^{-4}$, whereas the dedicated growth-rate wavelength ladder of Section~\ref{sec:validation_endtoend} uses a distinct $10^{-3}$ twelve-mode \emph{broadband} seed). Note that the magnetic mass-to-flux ratio $\mu_\Phi/\mu_{\Phi,\rm crit}$ is distinct from the thermal $f$; with $B_0$ set by $\beta$ and the central density by $f$ (at fixed $c_s$, background density, and $W_{\rm core}$), the on-axis mass-to-flux ratio scales as $\mu_\Phi/\mu_{\Phi,\rm crit}\propto f\sqrt{\beta}$, and every configuration in our grid has $\mu_\Phi/\mu_{\Phi,\rm crit}>1$, none is strictly magnetically subcritical; the low-$\beta$ runs are magnetically \emph{supported} (a strong field delays but does not prevent collapse).

The campaign's breadth maps the $(f,\beta,\mathcal{M},\theta)$ parameter space; the headline wavelength rests specifically on the growth-rate ladder of Section~\ref{sec:validation_endtoend}, and the large run count should not be read as statistical weight behind that value. We explored a base grid spanning Mach number $\mathcal{M}=0.5$--$5$ and plasma $\beta=0.05$--$5$ at a near-critical line mass, and a transition grid across $f=1.4$--$2.0$, $\beta=0.3$--$1.3$ and $\mathcal{M}=1$--$5$ to resolve the stable--unstable boundary. Full parameter ranges and run counts are given in Appendix~\ref{sec:campaigns} (Table~\ref{tab:campaigns}).

In brief (full detail in Appendix~\ref{sec:supplementary_campaigns_app}), the base grid establishes three outcome regimes: magnetically supported quasi-stability at strong fields ($\beta\lesssim0.15$), radial collapse, and longitudinal fragmentation; the near-critical, longitudinal-field case that is our focus beads at $t_{\rm frag}\approx1.1$--$1.2\,t_{\rm J}$, and turbulence changes the fragmentation \emph{wavelength} by $<5\%$ while shortening the timescale $3$--$5\times$. The radial-collapse/fragmentation timescale competition, the three-regime $(\beta,\mathcal{M})$ map, and the turbulence/critical-transition dependence are given in Appendix~\ref{sec:supplementary_campaigns_app}.


\subsection{Ambient confinement and the direct supercritical measurement}
\label{sec:modecompetition}

The outcome for a supercritical filament is set by whether it is confined by an ambient medium, as real HGBS filaments are within their parent clouds. Throughout this subsection $d$ denotes the transverse distance from the filament axis to the domain boundary, in units of $\lambda_{\rm J}$ (so $d=1.0$ places the boundary one Jeans length from the axis). The main 654-run grid uses a truncated Gaussian with no ambient medium, so the Gaussian skirt wraps through the periodic transverse boundary and feeds an accretion channel that suppresses longitudinal beading, so the filament collapses radially (Appendix~\ref{sec:regimes}). With ambient confinement (\texttt{filament\_ambient} generator, low-density ambient medium), the same configurations fragment longitudinally instead: a reconciliation suite (12 runs, $f\times\beta\times2$ seeds) beads at physical, $\beta$-dependent collapse times ($0.55$--$0.70\,t_{\rm J}$), converging across boundary conditions.

Ambient-confined supercritical filaments ($f=1.5$--$2.0$, $\theta=0^\circ$) fragment longitudinally: an ensemble of 39 runs (periodic and reflecting BCs, ambient confinement, $f\times\beta\times$multi-seed) develops 3--10 interior axial density maxima, at a spacing of order the Jeans length. We do not rely on the median peak spacing of this ensemble for a precise number, because counting peaks late in the collapse is epoch- and resolution-dependent (Section~\ref{sec:validation_endtoend}); we instead take the wavelength from the linear growth-rate spectrum, whose low-$n$ modes are resolution-stable and which gives an order-$\lambda_{\rm J}$, sub-classical value $\lambda_{\rm max}\approx1.0$--$1.14\,\lambda_{\rm J}$, $\approx0.35$--$0.5$ times the classical value (a factor $\approx2$--$3$; this is the peak of the converged part, not a resolved spectral maximum; Section~\ref{sec:validation_endtoend}). The scale declines mildly with both $f$ and $\beta$, so it is not constant across the supercritical regime. The ambient configuration confines the filament in a low-density medium (central-to-ambient density ratio $\sim10^2$) that is approximately in radial force balance, consistent with background-subtracted HGBS column-density profiles (crest-to-background contrasts of order $10$--$10^2$; \citealt{Arzoumanian2019}), so the confinement is representative of real filaments rather than tuned to produce fragmentation. The 39 runs span periodic ($n=12$, median $\lambda/W_{\rm fil}\approx0.89$), reflecting ($n=26$, median $\approx1.01$) and one user-outflow transverse boundary, and the periodic and reflecting distributions are statistically indistinguishable (KS $p=0.17$, though with these $n$ the test has limited power), consistent with combining them. This ensemble median ($\lambda/W_{\rm fil}\approx0.9$--$1.0$) and the growth-rate result ($\lambda_{\rm max}/W_{\rm FWHM}\approx2.0$, Section~\ref{sec:validation_endtoend}) are the same physical wavelength, $\lambda\approx1.1\,\lambda_{\rm J}$, expressed against different width normalisations (the broad beam-convolved Plummer $W_{\rm fil}$ versus the narrower Gaussian formation FWHM); they are not in tension. The ensemble beads early, at $t_{\rm bead}\approx0.6$--$0.75\,t_{\rm J}$, well before gravitational runaway and at moderate density contrast, so with $\lambda\approx\lambda_{\rm J}$ resolved by ${\sim}32$ cells it satisfies the Truelove criterion at the measurement epoch (Appendix~\ref{sec:validation}), unlike the perpendicular ambipolar runs below (which bead only at $C\approx250$--$560$, deep in collapse). The full per-run table ($f,\beta$, seed, BC, outcome, $t_{\rm bead}$, bead count, spacing) is released with the public data archive.

Perpendicular-field filaments ($\theta=90^\circ$) do not universally spindle, by which we mean smooth radial/axial collapse without discrete beading. An 18-run ideal-MHD grid at $d=1.0$ shows $7/18$ runs bead (at $\beta=0.5$ and $\beta=2.0$), while spindle collapse is concentrated near equipartition ($\beta\approx1$). This non-monotonicity reflects a competition between the axial fragmentation mode and the radial spindle mode: at $\beta\approx1$ the perpendicular field's tension response in the radial plane is maximised, so the radial mode reaches non-linearity before axial beading can develop, whereas at both weaker (thermally driven) and stronger (magnetic-pressure-delayed radial collapse) fields the axial mode has time to grow (a plausibility argument, not a mode-by-mode demonstration). Where perpendicular filaments do bead, the spacing appears short ($\lambda/W_{\rm fil}\approx0.4$, corrected conversion) --- but we caution that this \emph{conflicts} with linear theory, which for a static perpendicular field predicts a \emph{longer}, not shorter, mode; since the perpendicular runs are numerically unresolved ($N_J\approx1$--$2$, below the Truelove criterion), we cannot presently determine whether the short beading is physical or a resolution artefact, and draw no perpendicular conclusion in either direction. (A dedicated transverse-refinement ladder to $N_J\approx170$ now resolves this, Section~\ref{sec:optional_tests}: the ideal perpendicular filament spindles with no axial fragmentation, and the short beading is a Truelove artefact.) In an expanded 36-run non-ideal survey ($\eta_{\rm AD}=0.001$--$0.1$, $\beta=0.5$--$2.0$), $32/36$ ($89\%$) bead. We report this as suggestive only: because these runs are at $N_J\approx1$--$2$, where under-resolution is itself a known source of spurious fragmentation \citep{Truelove1997}, we cannot separate a genuine ambipolar enhancement from a numerical one. The median onset wavelength is $\lambda/W_{\rm fil}\approx0.6$ (bulk $0.6$--$0.7$; strongest case $\beta=0.5$, $\eta_{\rm AD}=0.01$ reaches $0.85$). These are fragmentation-onset measurements: the non-ideal runs bead at $t\approx0.44$--$0.55\,t_{\rm J}$ and undergo gravitational runaway essentially simultaneously (the density contrast is already $C\sim10^3$ at first beading and reaches $C\sim10^4$ within $\Delta t\approx0.02$--$0.04\,t_{\rm J}$, driving the CFL timestep below the kill threshold). A finer-cadence rerun ($\Delta t=0.02$) through the beading-to-collapse transition shows the measured wavelength does not reach a converged plateau over the available window. We therefore quote the median $\lambda/W_{\rm fil}\approx0.6$ explicitly as a fragmentation-onset value rather than a converged fragmentation wavelength. Ambipolar diffusion therefore promotes but does not yet resolve the perpendicular-field tension; a converged wavelength measurement is left to future work.

\subsection{Testing the resolution and magnetic-field alternatives}
\label{sec:optional_tests}
Three targeted campaigns test the concerns that most affect the interpretation: whether the longitudinal wavelength is resolution-limited, whether the perpendicular geometry produces genuine short beading, and whether a helical field --- the one configuration expected to shorten the axial mode --- actually does.

\emph{Resolution stability of $\lambda_{\rm max}$.} It is important to distinguish two refinement directions, because they behave oppositely and conflating them is the source of an apparent contradiction. Under \emph{transverse} refinement (the isotropic $128^3$--$512^3$ ladder, which improves the Jeans resolution of the collapsing spine) the short-wavelength part of the spectrum ($\lambda\lesssim0.9\,\lambda_{\rm J}$, modes $n\gtrsim9$) grows with resolution and is numerical (unresolved transverse sub-fragmentation; Fig.~\ref{fig:growthrate}a), whereas the physical band ($\lambda\gtrsim1\,\lambda_{\rm J}$, $n\leq8$) is converged. Under \emph{axial-only} refinement (holding the transverse grid fixed and doubling the axial resolution, $512\times128^2\to1024\times128^2$) the physical band reproduces to better than $0.5\%$ ($\Gamma$ at $n=6,7,8$, i.e.\ $\lambda=1.33,1.14,1.00\,\lambda_{\rm J}$, is $17.6,18.7,18.7$ at both resolutions --- rising to a plateau at the $\lambda\approx1\,\lambda_{\rm J}$ edge without turning over), \emph{and} the short-wavelength rise ($n\gtrsim9$) is likewise unchanged, confirming it is fixed by the axial seed-truncation/Nyquist scale (aliasing), not a physical mode that migrates to shorter $\lambda$ as the axial grid is refined. Both statements are consistent: the short-$\lambda$ excess is a transverse-under-resolution artefact that axial refinement neither creates nor removes. The $\lambda_{\rm max}\lesssim1\,\lambda_{\rm J}$ upper bound is therefore stable --- axial refinement does not move the physical-band plateau to shorter $\lambda$, and the only structure appearing at shorter $\lambda$ is the transverse-resolution artefact already identified.

\emph{The perpendicular geometry: a resolved spindle, not short beading.} On a transverse-refinement ladder ($128^3$--$512^3$) the near-critical \emph{ideal} perpendicular filament collapses monolithically in the radial direction at every resolution, with the dominant axial mode remaining at the box scale and the axial perturbation at the seed level throughout the well-resolved phase (at $512^3$ the spine Jeans length stays resolved with $N_J\approx170$ cells at contrast $C\approx3$, falling only to $N_J\approx40$ even at $C\approx30$ --- still $\approx10\times$ above the Truelove threshold; Appendix~\ref{sec:validation}). Short-wavelength axial structure ($\lambda/W\approx0.4$) appears only once the collapsing spine drops below $N_J\approx1$ cell per Jeans length; increasing resolution simply delays its onset to deeper collapse (Fig.~\ref{fig:perp_ladder}), the characteristic signature of a grid artefact. We therefore conclude that the perpendicular field \emph{suppresses} axial fragmentation in favour of radial collapse (anisotropic, flattening preferentially along the field rather than a true axisymmetric spindle) --- as static linear theory predicts (a perpendicular field lengthens the axial mode) --- and that the short beading reported at low resolution is numerical. This result is for \emph{ideal} MHD; the ambipolar-perpendicular case, where non-ideal diffusion may promote fragmentation (Section~\ref{sec:modecompetition}), remains separately unconverged, so the ideal-MHD spindle conclusion should be read with that caveat.

\emph{A seeded helical field pinches the width but does not shorten $\lambda_{\rm max}$; the force-balanced case remains untested.} We added a divergence-free helical field (uniform axial $B_z$ plus a toroidal $B_\phi(r)$ peaking at $r_c=W_{\rm core}$, built from a discrete curl of an axial vector potential $A_z(r)=B_t r_c\ln[1+(r/r_c)^2]$ so that $\nabla\!\cdot\!\mathbf{B}=0$ to machine precision) and scanned the toroidal-to-poloidal ratio $B_\phi/B_z=0$--$2$ ($B_\phi=0$ reproduces the longitudinal spectrum exactly, a code check). The axial growth-rate plateau peaks at $\lambda\approx1\,\lambda_{\rm J}$ for every toroidal fraction, with no excess growth at shorter wavelength even for a strongly-wound field ($B_\phi/B_z=2$; Fig.~\ref{fig:helical_scan}a). The toroidal component's actual effect is a radial \emph{pinch} that narrows the filament ($W_{\rm FWHM}$: $0.57\to0.43\,\lambda_{\rm J}$, $\approx25\%$; Fig.~\ref{fig:helical_scan}b) and accelerates radial collapse; it therefore raises $\lambda/W$ by reducing $W$, and so moves $\lambda/W$ \emph{towards} the classical value, not away from it. Crucially, this field is \emph{seeded}, not in force balance: the toroidal component is imposed on an already-contracting filament, so it drives an out-of-equilibrium radial pinch rather than establishing the confining toroidal equilibrium whose sausage ($m=0$) mode carries the shortened \citet{Fiege2000} wavelength. A seeded helical field therefore \emph{cannot} test that mechanism. We conclude that the static longitudinal and perpendicular geometries are disposed of as sources of a short axial spacing (longitudinal leaves $\lambda_{\rm max}$ unchanged; perpendicular spindles), but the helical/\citet{Fiege2000} pinch --- the one geometry linear theory predicts could shorten $\lambda$ --- remains \emph{untested}: it requires a force-balanced helical equilibrium (analogous to the drag-relaxed control of Appendix~\ref{sec:val_appendix}), which we defer to future work. Magnetic tension is thus not the demonstrated cause of the sub-classical scale, but neither is it excluded in its one relevant configuration.

\begin{figure}
\includegraphics[width=0.48\textwidth]{figures/fig_perp_ladder.pdf}
\caption{Ideal perpendicular-field refinement ladder ($128^3$--$512^3$; Section~\ref{sec:optional_tests}). \emph{(a)} Cells per local Jeans length $N_J$ at the beading-relevant spine contrast $C\approx3$: the collapsing spine is resolved far above the Truelove threshold ($N_J=4$, dashed) at every resolution. \emph{(b)} The spurious short-$\lambda$ axial ``beading'' emerges only once the spine drops below $N_J\approx1$, and refinement pushes that onset to ever-deeper collapse (at $512^3$ the filament stays monolithic to $N_J\approx40$ at $C\approx30$) --- the signature of a grid artefact, not a physical fragmentation mode.}
\label{fig:perp_ladder}
\end{figure}

\begin{figure}
\includegraphics[width=0.48\textwidth]{figures/fig_helical_scan.pdf}
\caption{Divergence-free helical-field scan ($B_\phi/B_z=0$--$2$, $512\times128^2$; Section~\ref{sec:optional_tests}). \emph{(a)} The axial growth-rate spectrum $\Gamma(\lambda)$ in the physical band is invariant to the toroidal fraction ($<6\%$, fit noise), with the plateau peak fixed at $\lambda\approx1\,\lambda_{\rm J}$; $B_\phi/B_z=0$ (black) reproduces the pure-longitudinal spectrum, a code check. No excess growth appears at shorter $\lambda$ for any winding. \emph{(b)} The one systematic effect is a radial \emph{pinch}: the formation FWHM narrows by $\approx25\%$, with onset at $B_\phi\gtrsim B_z$. The seeded helical field therefore raises $\lambda/W$ by shrinking $W$, not by shortening $\lambda$; the force-balanced \citet{Fiege2000} pinch mode is not tested here.}
\label{fig:helical_scan}
\end{figure}

\section{The fragmentation scale and its comparison with observation}
\label{sec:validation_endtoend}

The measured core spacing is sub-classical only when expressed against the filament FWHM. Because that comparison hinges on how the observable width is defined, and a single conversion factor was shown above to move the inferred $\lambda/W_{\rm fil}$ by a factor of two, we anchor the comparison on the quantity that carries no width-convention ambiguity, the spacing in units of the filament FWHM measured the same way on simulations and data, and we verify the whole chain end to end (Appendix~\ref{sec:val_appendix}).

First, the measurement is unbiased: passing synthetic filaments with a known spacing and width through the identical pipeline recovers $\lambda/W$ to better than $0.5\%$ over $\lambda/W=1.4$--$5.7$ (Appendix~\ref{sec:val_appendix}), confirming that the observed$\to$simulated comparison depends on the width normalisation convention rather than on any bias in the measurement itself. Second, forward-modelling the simulations end to end through the HGBS pipeline, projection to a column-density map, beam convolution, Plummer fitting and the same NN measurement, gives the observable $\lambda/W_{\rm fil}$ directly, with no analytic width conversion. The forward-model value ($\lambda/W_{\rm fil}\approx1.0$--$1.3$) is $\approx30\%$ above the analytic shortcut ($\approx0.9$, Eq.~\ref{eq:t1_apply}), because the shortcut applies a single $\sigma\!\to\!$FWHM factor whereas the forward model captures the full Gaussian$\to$beam-convolved-Plummer broadening; we adopt the forward-model value as definitive and carry this $\approx30\%$ offset as the width-conversion systematic in the comparison error budget.

The fragmentation wavelength must be measured with care. Counting density peaks late in the collapse gives an epoch- and resolution-dependent answer, because a collapsing filament continues to sub-fragment and higher resolution resolves finer peaks. We instead identify the dominant fragmentation mode from the growth-rate spectrum. Seeding the filament with a small broadband perturbation ($\delta v/c_s=10^{-3}$, twelve axial modes), we track the power $P(k,t)$ in each mode and fit its growth over the interval during which it grows approximately exponentially (Fig.~\ref{fig:growthcurves}). This growth accelerates as the filament collapses, so what we isolate is the dominant fragmentation mode of a \emph{dynamically collapsing} filament, not the mode of a quasi-static hydrostatic cylinder. The growth rate $\Gamma(\lambda)$ (Fig.~\ref{fig:growthrate}, left) rises as $\lambda$ decreases and is resolution-stable for $\lambda\gtrsim1\,\lambda_{\rm J}$ (the low-$n$ modes agree to $\sim5\%$ across $128^3$--$512^3$; $\Gamma$ at $n=7$ is $17.8$, $18.5$, $18.7$), rising to a \emph{plateau} at the short-wavelength edge of the converged band ($\Gamma$ at $n=6,7,8$, i.e.\ $\lambda=1.33,1.14,1.00\,\lambda_{\rm J}$, is $17.6,18.7,18.7$ --- flat within the fit noise across the last two modes), while below $\lambda\lesssim0.9\,\lambda_{\rm J}$ ($n\gtrsim9$) it separates and rises with \emph{transverse} resolution and is numerical. Because $\Gamma$ plateaus rather than turning over inside the converged band, we quote the dominant wavelength as an \emph{upper bound}, $\lambda_{\rm max}\lesssim1\,\lambda_{\rm J}$ (the plateau edge), and explicitly do \emph{not} claim a precisely-located spectral maximum (Appendix~\ref{sec:val_appendix}); the true dynamical peak could lie at shorter, transverse-unresolved wavelength. It is an order-$\lambda_{\rm J}$, sub-classical value: in formation-FWHM units $\lambda_{\rm max}/W_{\rm FWHM}\approx1.5$--$2.0$, i.e.\ a factor $\approx2.5$--$3$ below the classical $\lambda_{\rm m}\approx5$--$6\,W_{\rm FWHM}$ (\citealt{Fischera2012}), and at or below even the critical wavelength ($2.6\,W_{\rm FWHM}$) across $f=1.1$--$2.0$. (The classical $5$--$6$ uses the \emph{equilibrium} FWHM $\approx0.45\,\lambda_{\rm J}$; our contracting-run $W_{\rm FWHM}=0.58\,\lambda_{\rm J}$ gives $22H/W_{\rm FWHM}\approx4$, so we take the convention-free $\lambda_{\rm J}$-normalised statement $\lambda_{\rm max}\lesssim1.1$ vs $\lambda_{\rm m}=2.3\,\lambda_{\rm J}$ as primary.) The observations, however, use a beam-convolved Plummer width, so the \emph{like-for-like} value (below) is $\lambda/W\approx1.0$--$1.3$, not $\approx2$, and it is this matched value that should be set against the data (Table~\ref{tab:wave}).\footnote{The formation FWHM here is $W_{\rm FWHM}=0.58\,\lambda_{\rm J}$, the Gaussian FWHM fitted to the spine-averaged projected profile of these near-critical, ambient-confined runs at the growth-rate epoch. It is modestly narrower than the $W_{\rm form}=0.683\,\lambda_{\rm J}$ of the supercritical ensemble (Section~\ref{sec:width_correction}) because a near-critical filament has already begun to contract when its dominant mode is set; both are narrower than the initial Gaussian FWHM ($2.355\,W_{\rm core}=0.71\,\lambda_{\rm J}$) for the same reason. Table~\ref{tab:widthdefs} tabulates both formation widths. The classical fastest-growing wavelength is $\lambda_{\rm m}\approx5\,W_{\rm FWHM}$ (\citealt{Fischera2012}), so the like-for-like sub-classical statement is the density-independent ratio $\lambda_{\rm max}/W_{\rm FWHM}\approx2$ versus the classical $\approx5$--$6$, i.e.\ sub-classical by a factor $\approx2.5$--$3$ (and below even the classical critical wavelength $2.6\,W_{\rm FWHM}$).}

\begin{figure}
\includegraphics[width=0.47\textwidth]{figures/fig_growthcurves.pdf}
\caption{Power in three axial modes versus time for a near-critical filament ($f=1.1$). Each grows with an approximately exponential interval (the fit window, shaded) whose slope gives the growth rate; the growth steepens as the filament collapses, so the extracted rate characterises the dominant mode of a dynamically collapsing filament. The peak mode $n=7$ grows faster than the long-wavelength $n=3$, giving the sub-classical $\lambda_{\rm max}$; the grid-scale $n=12$ is shown for reference.}
\label{fig:growthcurves}
\end{figure}

\begin{figure*}
\includegraphics[width=0.92\textwidth]{figures/fig_growthrate.pdf}
\caption{\emph{(a)} Growth-rate spectrum $\Gamma(\lambda)$ of the axial fragmentation modes for a near-critical filament ($f=1.1$) at three resolutions. For $\lambda\gtrsim1\,\lambda_{\rm J}$ the three curves overlap: the growth rate is resolution-converged. At the grid scale ($\lambda\lesssim0.9\,\lambda_{\rm J}$, shaded) the curves separate and rise with resolution, the numerical signature of unresolved sub-fragmentation. The low-$n$ modes are resolution-stable and $\Gamma$ plateaus (does not turn over) at the short-wavelength edge of the converged band, so $\lambda_{\rm max}\lesssim1\,\lambda_{\rm J}$ is an \emph{upper bound} (the plateau edge), not a resolved maximum (Appendix~\ref{sec:val_appendix}). \emph{(b)} The resulting $\lambda_{\rm max}/W_{\rm FWHM}$ versus line-mass fraction (Gaussian FWHM; points: individual runs; dashed: median), well below the classical $\lambda_{\rm m}/{\rm FWHM}\approx5$--$6$ at every $f$. The like-for-like comparison with the observations uses a beam-convolved Plummer width $\sim1.5\times$ broader, giving matched $\lambda/W\approx1.0$--$1.3$ (grey band): this, not the Gaussian $\approx2$ (upper limit), is what should be set against the observed $\lambda/W\approx0.4$--$1.9$.}
\label{fig:growthrate}
\end{figure*}

The convention-independent statement is that the dominant fragmentation wavelength of a dynamically collapsing filament is sub-classical. In the cleanest, width-free form it is $\lambda_{\rm max}\lesssim1.0$--$1.14\,\lambda_{\rm J}$ versus the classical $\lambda_{\rm m}=22H=2.3\,\lambda_{\rm J}$, a factor $\approx2$; we take this $\lambda_{\rm J}$-normalised statement as primary. In FWHM units it is $\lambda_{\rm max}/W_{\rm FWHM}\approx2$ versus the classical $5$--$6$ (\citealt{Fischera2012}; $\lambda_{\rm m}=4\times$ the flat diameter, \citealt{Inutsuka1992}), a factor $\approx2.5$--$3$ --- slightly larger only because the contracting-run FWHM ($0.58\,\lambda_{\rm J}$) exceeds the equilibrium FWHM ($\approx0.45\,\lambda_{\rm J}$) on which the classical $5$--$6$ is defined --- and at or below even the critical wavelength ($2.6\,W_{\rm FWHM}$). This is an \emph{upper bound}: the spectrum rises to a plateau at the short-wavelength edge of the resolution-converged band ($\lambda\approx1\,\lambda_{\rm J}$) and does not turn over there, so the true peak may lie at shorter, transverse-unresolved wavelength (the ratio $\lambda_{\rm max}/\lambda_{\rm m}$ measured this way runs $\approx0.35$--$0.5$ across $f=1.1$--$2.0$ (a factor $\approx2$--$3$), but each value is an upper bound). The selected axial \emph{mode} is imprinted early in its own growth and is fit-window-invariant (the same mode dominates in every linear sub-window; Appendix~\ref{sec:val_appendix}), although by that epoch the filament has already contracted radially to $C\approx20$--$40$ (where the Jeans length is still resolved, $N_J\gtrsim5$; Appendix~\ref{sec:validation}). We therefore quote the \emph{density-independent} ratio $\lambda/W_{\rm FWHM}$, with both lengths co-measured at the same epoch; the equilibrium-recovery control (Appendix~\ref{sec:val_appendix}) confirms this ratio returns the classical $5$--$6$ for a true equilibrium and $\approx1.4$--$2$ for a contracting one, so the sub-classical ratio is physical and not an artefact of the density at which the classical scale is evaluated. The shortening is the direction predicted for the accreting/contracting filament of \citet{Clarke2016,Clarke2017}. Two caveats temper the causal attribution. (i) The fitted growth rates ($\Gamma\approx18$--$26$ in units of $t_{\rm J}^{-1}$) exceed the free-fall rate ($\Gamma_{\rm ff}=\sqrt{4\pi G\rho_0}\approx2\pi$), i.e.\ they are not quasi-static linear eigenvalues but the accelerating growth of a collapsing background (the fitted $\Gamma$ correspond to the local free-fall rate at an effective density $\rho_{\rm eff}/\rho_0\approx(\Gamma/2\pi)^2\approx8$, distinct from the peak central contrast $C\approx20$--$40$; we use only the mode \emph{ordering}, not the absolute rate); the \emph{wavelength} (peak of $\Gamma(\lambda)$) is what we use, and it is corroborated by the fit-window invariance and the nonlinear bead spacing, not the absolute rate. Crucially, the selected wavelength is \emph{not} the trivial instantaneous Jeans scale of the contracted core: at the measurement epoch the on-axis density is $\rho_{\rm eff}\approx8\rho_0$, so the instantaneous Jeans length is $\lambda_J(\rho_{\rm eff})\approx\lambda_J(\rho_0)/\sqrt8\approx0.35\,\lambda_J(\rho_0)$, yet the measured $\lambda_{\rm max}\approx1\,\lambda_J(\rho_0)\approx2.8\,\lambda_J(\rho_{\rm eff})$ --- nearly three times the instantaneous Jeans length. The mode is therefore imprinted \emph{earlier}, when the density was lower, and does \emph{not} track $\lambda_J(\rho_c(t))$ as the core contracts (consistent with its fit-window invariance); were $\lambda_{\rm max}$ simply following the shrinking Jeans scale it would sit near $0.35\,\lambda_J(\rho_0)$, well inside the transverse-unresolved numerical band, which it does not. It is, however, close to the \emph{classical} fastest-growing wavelength evaluated at the elevated central density, $22H(\rho_{\rm eff})\approx22H(\rho_0)/\sqrt8\approx0.8\,\lambda_J(\rho_0)$: the filament fragments at approximately the classical scale \emph{for its current density}. This is the crux of the physical interpretation, and we state it plainly: the sub-classical \emph{deficit} in $\lambda/W$ arises not because the fragmentation physics is non-classical, but because the observable FWHM $W$ does \emph{not} contract in step with $\rho_c$ (the extended envelope lags the collapsing spine), so $\lambda\sim22H(\rho_c)$ shrinks while $W$ stays broad and $\lambda/W$ falls below the equilibrium $5$--$6$, in which $W$ and $22H$ are locked together. The equilibrium-recovery control confirms exactly this: when $W$ and $22H$ \emph{are} locked (a genuine equilibrium) the identical pipeline returns $\lambda/W\approx5$--$6$. (The two density markers refer to different baselines --- the box-mean contrast $C\equiv\rho_{\rm max}/\langle\rho\rangle_{\rm box}\approx20$--$40$ versus the on-axis $\rho_{\rm eff}/\rho_0\approx8$ that sets $\Gamma\approx18$ --- and are consistent statements about the same evolution differing by the box-filling factor; the Truelove bookkeeping uses the on-axis density and gives $N_J\gtrsim5$ at measurement.) (ii) Our initial profile is Plummer-like ($p\simeq2$), matching observed filaments but not the isothermal equilibrium ($p=4$), so it contracts from the start; a $p=4$ Ostriker profile gives the same order-$\lambda_{\rm J}$ wavelength ($\approx0.8\,\lambda_{\rm J}$), so the deficit is not a profile-index artefact, but that cylinder is itself only approximately relaxed, so what we demonstrate is the generic ``fragmentation from a contracting state'' rather than specifically the accretion-driven Clarke mechanism (Section~\ref{sec:limitations}).

The simulated observable is quoted at several values in this paper, and it is worth converting once. The convention-independent internal ratio is $\lambda_{\rm max}/W_{\rm core}\approx3.3$; the $\sigma$-normalised shortcut ($\times0.27$; Eq.~\ref{eq:t1_apply}) gives $\lambda/W_{\rm fil}\approx0.9$, whereas forward-modelling the same runs through the HGBS pipeline (projection, beam convolution, Plummer fit) gives $\lambda/W_{\rm fil}\approx1.0$--$1.3$. We adopt the forward-model value as \emph{the} observable, since the shortcut normalises by the initial-profile $\sigma$ rather than the evolved beam-convolved Plummer FWHM; the $0.9$ figure is a shortcut lower estimate only. In Gaussian formation-FWHM units the same wavelength is $\lambda_{\rm max}/W_{\rm FWHM}\approx1.5$--$2.0$, but a Gaussian FWHM is $\sim1.5$--$1.9\times$ narrower than the observed Plummer FWHM, so the ``$\approx2\approx1.9$'' coincidence is not like-for-like. In matched units the simulated $\lambda/W\approx1.0$--$1.3$ is the value to set against the data (at the order-of-magnitude level; the observed spacing is a line-density statistic, Section~\ref{sec:clustering}); its convention-by-convention standing is given below and in Section~\ref{sec:error_budget}. We stress this comparison is \emph{resolution-limited as well as convention-limited}: the simulated wavelength is an upper limit that could shift to shorter $\lambda$ at higher resolution, which would move single-filament fragmentation from sufficiency toward insufficiency.

Single-filament gravitational fragmentation, evolved dynamically, thus produces a converged \emph{upper bound} on the fragmentation wavelength that is sub-classical and consistent with the observed spacing for the near-critical filaments that dominate the sample. We are careful here about what is converged: the growth-rate \emph{spectrum} does not turn over within the resolution-converged band, so what is bounded is the wavelength ($\lambda_{\rm max}\lesssim1\,\lambda_{\rm J}$), not a located maximum; the object measured is the dominant mode of a \emph{dynamically collapsing} configuration, for which the linear-eigenmode concept is only approximate, so the robust deliverable is the density-independent \emph{ratio} $\lambda/W$ together with its equilibrium-recovery calibration, not a precise linear ``$\lambda_{\rm max}$''. What secures the sub-classical \emph{sign} as physical is the drag-relaxed equilibrium-recovery control (Appendix~\ref{sec:val_appendix}), in which the identical pipeline recovers the classical turnover at $22H$, whereas the contracting runs have no turnover down to the grid-limited band. We stress that this is a numerical \emph{instance}, in a controlled MHD campaign, of the generic sub-hydrostatic expectation that a still-contracting (non-equilibrium) filament fragments below the marginally-stable hydrostatic mode \citep{Clarke2016,Clarke2017,Clarke2020} --- because $\lambda_{\rm J}\propto\rho_c^{-1/2}$ shrinks as the central density rises --- not the discovery of a new mechanism; our distinct contributions are the numerical demonstration for the longitudinal near-critical case, the equilibrium-recovery control, and the homogeneous Gaia-DR3 reanalysis. The observations are \emph{consistent with} ordinary single-filament fragmentation under the DisPerSE convention (Section~\ref{sec:error_budget}), but --- because the simulated scale is only an upper bound --- this establishes consistency, not sufficiency: a shorter true wavelength would make single-filament fragmentation \emph{over}-produce cores. Neither uniqueness nor sufficiency is demonstrated, and hierarchical-fibre projection (Section~\ref{sec:fibre_projection}) and magnetic tension remain viable but, within the factor-$\sim$2 systematics, are neither demanded nor excluded. The central sub-classical claim is falsifiable, and we state the criterion explicitly: it is \emph{refuted} if, on resolved single filaments (or individual velocity-coherent fibres) measured with a fixed width convention, the core line-density is no higher than the classical one-core-per-$\lambda_{\rm m}$ (i.e.\ cores spaced at $\lambda_{\rm m}\approx5$--$6\times$FWHM, as sparse as the equilibrium instability predicts). Two further discriminants follow: fibre-resolved kinematics placing cores at the classical spacing \emph{within} individual fibres would favour projection over single-filament fragmentation, and more strongly supercritical filaments should show shorter core spacing (Section~\ref{sec:limitations}).

\section{DISCUSSION}
\label{sec:discussion}

\subsection{Why the spacing matters: fragment masses and the CMF}
\label{sec:cmf}
The spacing matters because, with the line mass, it sets the characteristic mass of the fragments that become prestellar cores and so seed the core mass function (CMF) and ultimately the stellar IMF. For a filament of line mass $M_{\rm line}$ fragmenting at spacing $\lambda$, the mean fragment mass is $M_{\rm frag}\sim\epsilon\,M_{\rm line}\,\lambda$ ($\epsilon\lesssim1$ the reservoir fraction). A near-critical $M_{\rm line}\approx16\,M_\odot\,{\rm pc^{-1}}$ ($T\approx10$\,K) and $\lambda\approx0.1$--$0.2$\,pc give $M_{\rm frag}\sim1.6$--$3\,M_\odot$, i.e.\ $\sim0.5$--$1\,M_\odot$ for $\epsilon\sim0.3$, consistent with the HGBS CMF peak. The \emph{sub-classical} spacing is what makes this work: the classical $\approx5$--$6\times$FWHM would place cores $3$--$4\times$ farther apart and predict $3$--$4\times$ more massive fragments, overshooting the CMF peak. We present this as an order-of-magnitude \emph{consistency check}, not a fit: $\epsilon$ and $M_{\rm line}$ are degenerate (only their product $\epsilon M_{\rm line}$ enters, and we have not measured $\epsilon$ independently), and because the observable is a line density rather than a resonant wavelength (Section~\ref{sec:clustering}) this predicts a \emph{mean} fragment mass, not a peaked CMF. The robust point is the sign --- sub-classical spacing lowers characteristic masses toward the observed range --- not a numerical CMF prediction.


\subsection{Hierarchical-fibre projection is consistent with the sub-classical spacing}
\label{sec:fibre_projection}

Section~\ref{sec:validation_endtoend} showed that single-filament dynamical fragmentation is consistent, within present methodological uncertainties, with the observed sub-classical spacing. A second, purely geometric route reaches the same value without any modified physics, and is worth setting out because the substructure it relies on is independently observed: high-resolution kinematic studies show that HGBS filaments are not monolithic cylinders but bundles of $N$ velocity-coherent \emph{fibres} \citep{Hacar2013,Hacar2018,Yang2024}. This motivates a purely geometric alternative to modified fragmentation physics: if each fibre fragments near the classical wavelength but core spacing is measured \emph{across} the bundle, as any filament-level skeleton pipeline necessarily does, the interleaving of cores from $N$ fibres compresses the apparent filament-level spacing. The premise, quasi-periodic fibre-to-core fragmentation at the classical scale, is a hypothesis rather than an established measurement (\citealt{Yang2024} find only a weak such imprint), so we present projection as viable but unconfirmed.

We tested this quantitatively with a forward-model that reuses our exact NN measurement (order cores along the reconstructed spine by arclength; take adjacent separations; median $=\lambda$). Each fibre is a cylinder of physical width $W_{\rm fib}$ that fragments at its \emph{own} classical wavelength $\lambda_{\rm fib}\approx5$--$6\,W_{\rm fib}$ (\citealt{Fischera2012}; the FWHM-based fastest-growing scale). We populate a filament with $N$ such fibres (random inter-fibre phase, $\sim30\%$ lognormal intrinsic scatter), project the cores onto the shared spine, and measure the spacing in units of the \emph{filament} width $W_{\rm fil}$. This assumes fibres are quasi-hydrostatic and fragment at the classical scale, which sits uneasily with this paper's own result. The physically consistent cases are limited: a genuinely \emph{subcritical} fibre is gravitationally stable and would not fragment into cores at all, whereas a \emph{near-critical} one, being self-gravitating and contracting, would by our own argument fragment sub-classically at $\approx\tfrac12\lambda_{\rm fib}$. In the latter case the numerator halves and the $r/N$ needed to match the data \emph{doubles} (Eq.~\ref{eq:fourrn} becomes $\lambda/W_{\rm fil}\approx(2.5$--$3)r/N$). We therefore treat the classical-fibre $(5$--$6)r/N$ result as an upper bound on the projection effect; the fibre and single-filament pictures are not independent. Two distinct hierarchical effects then compress the apparent filament-level spacing below the single-cylinder value: the fibres are narrower than the filament ($r\equiv W_{\rm fib}/W_{\rm fil}<1$), and there are several of them. In the interleaving limit the Monte-Carlo model reproduces the simple scaling
\begin{equation}
\frac{\lambda}{W_{\rm fil}} \;\approx\; \frac{(5\text{--}6)\,r}{N},
\label{eq:fourrn}
\end{equation}
verified across $r$ and $N$ in Fig.~\ref{fig:fibre_degen}. Either effect alone can produce a sub-classical value: a single fibre of half the filament width ($r=0.5$, $N=1$) gives $\lambda/W_{\rm fil}=2.0$, as do two full-width fibres ($r=1$, $N=2$).

The observation constrains only the combination $r/N$, not $N$ individually. The fixed-width value $\lambda/W_{\rm fil}\approx1.9$ (periodic-corrected; Table~\ref{tab:headline}) requires $r/N\approx0.35$; the region-specific Plummer value ($\approx1.4$, raw) gives $r/N\approx0.25$; and the independent FilFinder skeletons ($\lambda/W\approx0.4$--$0.8$) give $r/N\approx0.07$--$0.15$. The observed fibre widths, $W_{\rm fib}\approx0.03$--$0.05$\,pc (so $r\approx0.3$--$0.5$; \citealt{Hacar2013,Arzoumanian2011}), then imply a small fibre number, $N\approx1$--$2$ for the DisPerSE-based spacings and $N\approx3$--$5$ if the FilFinder spacings are correct. \citet{Hacar2013} resolve of order $3$--$5$ velocity-coherent fibres per bundle in Taurus, which with $r\approx0.3$--$0.5$ gives $r/N\approx0.07$--$0.15$, below the $\approx0.35$ that the DisPerSE spacing requires. Taken at face value, the projection model with the observed fibre counts would \emph{over}-compress the spacing (predict cores closer than observed), matching the DisPerSE value only at the low end $N\approx1$--$2$; it is internally consistent only for a restricted region of $(r,N)$, a further reason it is not uniquely favoured. We do \emph{not} claim a specific fibre number; the robust statement is that measuring core separations across hierarchical structure (narrower, and generally multiple, fibres) in units of the filament width reduces the apparent spacing below the classical prediction by the observed factor, with no modified fragmentation physics required.

The $(5$--$6)r/N$ scaling is also reproduced when the forward-model is run on the real DisPerSE arclengths and counts of the four robust regions rather than idealised straight spines; that exercise is only a consistency check, since it seeds classical-spacing cores onto the observed skeletons and so assumes the fibre population rather than measuring it.

\begin{figure}
\includegraphics[width=0.48\textwidth]{figures/fig_fibre_degeneracy.pdf}
\caption{Hierarchical-fibre forward-model, showing the filament-level $\lambda/W_{\rm fil}$ recovered by our NN pipeline as a function of the fibre-to-filament width ratio $r=W_{\rm fib}/W_{\rm fil}$ and the number of fibres $N$; the model follows $\lambda/W_{\rm fil}\approx(5$--$6)r/N$ (Eq.~\ref{eq:fourrn}). White contours mark the observed fixed-width and region-specific values, the orange contour the FilFinder value, and the red band the observed fibre-width range ($r\approx0.3$--$0.5$). The observation constrains the combination $r/N$, not $N$ alone: the sub-classical value is reproduced by a single narrow fibre, by a few wider fibres, or by intermediate combinations.}
\label{fig:fibre_degen}
\end{figure}

Hierarchical structure, the combination of narrower fibres and interleaving that compresses the apparent spacing below the classical filament-level value (Eq.~\ref{eq:fourrn}), is thus a natural and sufficient way to produce the observed sub-classical spacing without modified magnetic physics, and it is the natural reading of the fibre-resolved observations \citep{Yang2024,Hacar2013}; it also accounts for the large skeleton-algorithm and $L/3$ systematics ($\S$\ref{sec:error_budget}), as expected if the measurement slices across blended structure. It is not, however, \emph{required}: the end-to-end validation (Section~\ref{sec:validation_endtoend}) shows single-filament dynamical fragmentation already accounts for the bulk of the spacing, so projection is one of at least two viable readings, and it assumes both the fibre population and quasi-periodic fibre-level fragmentation rather than measuring them. It does not \emph{disprove} a magnetic contribution, only that one is not required and that present data cannot distinguish the possibilities. The decisive test is to decompose the observed cores into their parent velocity-coherent fibres and measure the fibre-level spacing across the robust regions (Section~\ref{sec:limitations}): a quasi-periodic classical-scale spacing within individual fibres, with the fibre widths, would break the $r/N$ degeneracy and establish the hierarchical origin. A complementary theoretical test, not attempted here, is a self-consistent MHD (or SPH) fragmentation calculation of a multi-fibre bundle run through the same NN pipeline. Numerical work reaches the same under-determined conclusion from the theory side: \citet{Clarke2020} find that fragmentation proceeding through sub-filaments leaves \emph{no} single characteristic core-spacing length-scale, consistent with our finding that the measured spacing is convention-limited rather than a universal wavelength.

A related coherence point deserves statement: our single-filament simulations fragment into a \emph{quasi-periodic} comb of beads at $\lambda_{\rm max}$, whereas the observed cores are random-to-clustered with no comb (Section~\ref{sec:clustering}). The two are reconciled only if sub-fragmentation, interleaving of multiple fibres, and projection across blended structure erase the periodicity and broaden the comb into the observed near-exponential gap distribution --- as \citet{Clarke2020} find for sub-filament fragmentation. We do not simulate this multi-fibre stage; a self-consistent MHD/SPH fragmentation of a fibre bundle run through the same NN pipeline is the missing link, and until it is done the match between a single-filament $\lambda_{\rm max}$ and the observed line density is an order-of-magnitude, not a morphological, agreement.



\subsection{Field geometry, boundary conditions and the perpendicular case}
\label{sec:why_mhd_fails}

Across the campaign the \emph{outcome} of a run, radial collapse or longitudinal beading, depends on regime and boundary conditions: an ambient-confined supercritical filament fragments longitudinally, whereas the identical filament in a periodic box collapses radially. A single filament therefore has no setup-independent fate. Its dominant fragmentation \emph{mode} is nonetheless well defined and resolution-stable (Section~\ref{sec:validation_endtoend}), giving a converged \emph{upper bound} $\lesssim0.5$ times the classical value, sub-classical and consistent with the observed spacing. The perpendicular-field case is the exception: it is numerically unconverged ($N_J\approx1$--$2$ cells per Jeans length; Appendix~\ref{sec:validation}) and places no quantitative constraint. It is sometimes argued that most filaments are threaded by perpendicular fields, from the tendency of dense filaments to lie perpendicular to the large-scale field seen by Planck \citep{Planck2016}; the projected large-scale orientation need not equal the three-dimensional internal field geometry of a dense filament, and hourglass morphologies in which the field turns longitudinal inside the crest are observed \citep{Jadhav2026}. A resolved perpendicular filament spindles rather than fragmenting axially (Section~\ref{sec:optional_tests}), so it is not a source of tension: perpendicular filaments do not produce the observed axial core chains.

Where perpendicular geometry does apply, three non-exclusive resolutions of the residual deficit remain: hierarchical fibre fragmentation (bundles of near-critical fibres, each fragmenting longitudinally and projected across the bundle; \citealt{Hacar2013,Hacar2018,Yang2024}); non-ideal MHD at later evolution, which narrows the factor-$\sim3$ perpendicular deficit in longer runs; and an observational lag, in which the observed cores are relics of an earlier fragmentation episode. The two leading interpretations, single-filament dynamical fragmentation and hierarchical projection, are distinguished decisively by fibre-resolved core spacing across the robust regions and by high-resolution polarimetry of filament interiors (Section~\ref{sec:limitations}); the temperature-gradient signature (Section~\ref{sec:mt_variation}) adds a third test, predicting spatially differential fragmentation that fibre geometry alone would not produce.

Sub-isothermal equations of state ($\gamma<1$) accelerate fragmentation by $30$--$40\%$ without changing the wavelength, while adiabatic heating ($\gamma=5/3$) suppresses it entirely (0\% fragmentation to $t>30$--$40\,t_{\rm J}$; Appendix~\ref{sec:validation}); the isothermal suite is therefore the conservative (most unstable) case in a globally averaged sense. This framing does not apply locally: Section~\ref{sec:mt_variation} shows the core dust temperature, and hence the local Jeans length, varies spatially by $\sim12\%$ within active clouds, driven by a spatially varying radiation field, so sub-regions can fragment at a scale that differs from the isothermal prediction in either direction. This is a directly measured effect, since an imposed temperature gradient of this magnitude shifts the local $\lambda/W$ by $\sim0.1$ ($\sim10$\,per-cent) and relocates where cores form (Section~\ref{sec:mt_variation}).

\subsubsection{Non-Ideal MHD Effects}
\label{sec:nonideal}

At the densities probed here ($n_0\sim10^4$\,cm$^{-3}$), the ambipolar-diffusion time ($t_{\rm AD}\propto L^2/\eta_{\rm AD}$; $\sim10$--$20\,t_{\rm ff}$ for typical filament conditions) is long compared with the collapse time \citep{Tassis2004}, so ideal MHD is a good approximation for the longitudinal fragmentation stage: four non-ideal runs at $\eta_{\rm AD}=\{0.01,0.05\}$ still collapse radially, $\sim12$--$15\%$ \textit{earlier} than the ideal control (flux leakage erodes magnetic support, the conservative direction). Non-ideal effects nonetheless change the outcome for perpendicular fields, where ambipolar diffusion promotes longitudinal fragmentation (Section~\ref{sec:modecompetition}); Hall and Ohmic terms remain relevant for late-stage core evolution.

\subsection{Limitations and future work}
\label{sec:limitations}

Two limitations bound the strength of our conclusion. First, on the observational side, the characteristic $\lambda/W$ and its cloud-to-cloud scatter ($\sigma\approx0.4$) rest on only four independent regions (Orion~B, Aquila, Perseus, Taurus). These are all nearby, low-mass Gould Belt clouds of broadly solar metallicity and similar interstellar radiation field, so our result is established only for that regime; whether the sub-classical spacing persists in high-mass star-forming complexes, at lower metallicity, or under stronger irradiation (all of which alter the thermal Jeans scale and the critical line mass) is untested and is a natural target for future surveys. Four clouds are also a thin base for an intrinsic-scatter estimate or for claiming that the primary $\lambda/W\approx1.4$ (region-Plummer) is representative of filaments generally, and the value itself carries the factor-of-two skeleton-algorithm and width-convention systematics quantified in Section~\ref{sec:error_budget}. Extending the full-coverage nearest-neighbour pipeline to further HGBS regions, and to independent surveys (e.g.\ JCMT/SCUBA-2 Gould Belt fields), is needed to confirm that neither the central value nor the scatter is dominated by these four clouds. Second, on the theoretical side, the simulations are deliberately idealised controlled experiments, not models of specific filaments. The periodic-BC main grid does not represent the turbulent, pressure-supported, magnetically coupled boundaries of a real filament embedded in its parent cloud; our ambient-confined configuration is an idealisation of that confinement. Neither turbulent inflow/outflow at the boundaries nor gravitational coupling to the parent cloud is modelled, so the simulations are controlled experiments rather than direct models of individual HGBS filaments.

On the causal attribution (expanded in Section~\ref{sec:validation_endtoend}): our seeded filaments contract from the start, so what we demonstrate is the general ``fragmentation from a non-equilibrium, contracting state'' rather than specifically the accretion-driven case of \citet{Clarke2016,Clarke2017}; a near-equilibrium filament grown by controlled accretion is the decisive future test. The simulations adopt several further idealisations: (i) an \emph{isothermal} equation of state (sub-isothermal cooling accelerates fragmentation without changing the wavelength, while the temperature-gradient run of Section~\ref{sec:mt_variation} shifts the local $\lambda/W$ by $\sim10$\,per-cent, so the isothermal values are locally uncertain at the $\sim10$--$20\%$ level); (ii) \emph{periodic transverse boundaries}, whose accretion channel suppresses fragmentation unless ambient confinement is imposed (Section~\ref{sec:modecompetition}); (iii) \emph{simplified initial profiles} that omit the sub-structure of real filaments; and (iv) a \emph{prescribed, low-amplitude Kolmogorov turbulence spectrum} rather than driven transonic turbulence, which leaves the wavelength unchanged ($<5\%$) but shortens the timescale by $3$--$5\times$ (Appendix~\ref{sec:supplementary_campaigns}).

The fragmentation \emph{classification} is resolution-independent (the collapse \emph{time} itself carries an $8.5\%$ offset between $128^3$ and $256^3$ and is not claimed converged; Appendix~\ref{sec:validation}) and domain-doubling changes the collapse time by $<2\%$, but the transition zone is sampled with only 2 seeds per point. Non-linear core merging could in principle inflate the apparent supercritical spacing by up to $\sim50\%$ over a few Myr, but we have verified it does \emph{not} bias the quoted value: $\lambda/W_{\rm fil}$ is measured at maximum bead count (fragmentation onset), and in $38/39$ runs the peak count is still rising at the final snapshot while no run shows the peak-count decline that would signal merging. The ambipolar wavelengths are likewise fragmentation-onset measurements (the runs reach gravitational runaway before a converged plateau is established; Section~\ref{sec:modecompetition}).

The decisive follow-ups are: (i) fibre-resolved NN spacing on the full HGBS catalogues, to test the hierarchical interpretation; (ii) high-resolution polarimetry of filament interiors, to test the longitudinal-versus-perpendicular assumption; (iii) MHD simulations with a spatially varying temperature field (Section~\ref{sec:mt_variation}); (iv) longer non-ideal runs, to test whether the perpendicular deficit closes; (v) a seeding-fixed domain study to \emph{refine} (rather than bound) the near-critical $\lambda_{\rm max}$, the only remaining lever on its precision (the seed-realisation and open-gravity-boundary robustness having been established, Appendix~\ref{sec:val_appendix}).


\subsection{Non-isothermal effects and the temperature-gradient test}
\label{sec:mt_variation}

The MHD suite is isothermal, whereas HGBS dust-temperature maps vary by $\sim12$\,per-cent ($\sim10$--$13$\,K), a local perturbation since $\lambda_{\rm J}\propto\sqrt{T}$. A single illustrative run with an imposed $\sim12$\,per-cent longitudinal temperature gradient fragments \emph{differentially} (the cool, locally supercritical half beads while the warm half is suppressed), shifting the local wavelength by $\sim10$\,per-cent, because temperature sets both the Jeans scale and the critical line mass ($\mu_{\rm crit}\propto T$). The isothermal assumption is thus conservative only in a globally averaged sense; a self-consistent treatment is deferred to future work.


\section{CONCLUSIONS}

Our primary results are methodological, and we state them first because they are what the data robustly support: (i) the observed core spacing is sub-classical in \emph{sign} under every skeleton/width convention, but its \emph{magnitude} is dominated by measurement convention (skeleton algorithm, width definition, injection--recovery model), which shifts $\lambda/W$ by nearly a factor of two --- larger than any physical signal we can presently isolate; (ii) the nearest-neighbour spacing is the appropriate observable, whereas the widely-used pairwise median measures the region extent $L/3$; and (iii) a dynamically contracting filament \emph{naturally} produces a sub-classical fragmentation scale, validated against a hydrostatic equilibrium control. We do \emph{not} claim to have identified the physical cause of the deficit. Current observations cannot uniquely discriminate single-filament dynamical fragmentation from hierarchical fibre substructure, the two mechanisms that remain viable (below) --- and which, in the self-consistent near-critical case, may not even be physically distinct (Section~\ref{sec:fibre_projection}). The residual uncertainty is dominated by observational convention; the simulated wavelength, while resolution-stable in the converged band (Section~\ref{sec:optional_tests}), is an \emph{upper bound} whose true value may be shorter (the growth-rate spectrum does not turn over), so what the simulations establish is the sub-classical \emph{sign}, corroborated by the equilibrium-recovery control, not a precise wavelength. We address the magnetic geometries directly (Section~\ref{sec:optional_tests}): an \emph{ideal-MHD} perpendicular field, resolved to $N_J\approx40$--$170$ (over spine contrast $C\approx30$--$3$), drives monolithic radial collapse with no axial fragmentation --- so it does not produce short beading, and a static perpendicular field of the kind Planck indicates is common cannot be the source of the sub-classical scale; whether ambipolar diffusion changes this remains numerically unresolved (our unconverged non-ideal runs bead in $32/36$ cases, Section~\ref{sec:modecompetition}). The one geometry that linear theory predicts \emph{could} shorten the wavelength is helical/toroidal; in our seeded (non-equilibrium) runs it narrows the width by a radial pinch rather than shortening $\lambda$, but because the destabilising \citet{Fiege2000} pinch operates on a force-balanced toroidal equilibrium --- which we do not simulate --- magnetic tension remains \emph{untested} in its one relevant configuration, rather than demonstrated or excluded.

We have investigated the origin of the sub-classical core spacing in HGBS filaments, combining an observational re-analysis, a large MHD simulation campaign, and a hierarchical-fibre forward-model. Our main conclusions are:

\begin{itemize}
    \item \textbf{Established.}
    \begin{itemize}
        \item Gaia DR3 distances raise the core-weighted mean spacing modestly and region-dependently ($\sim10\%$ for the robust regions; large for Aquila/Serpens, negligible for the dominant Orion~B, whose pre-Gaia HGBS distance was already $\sim400$\,pc). The absolute $\lambda/W$ uses the adopted Gaia distances and does not depend on this revision.
        \item The primary physical core spacing is $\lambda/W\approx1.4$ (median of direct region-specific Plummer fits); the fixed $0.10$\,pc-width value is $\lambda/W\approx1.2$ under the data-matched clustered injection--recovery correction (the cores are random-to-clustered, not periodic; Section~\ref{sec:clustering}), rising to $\approx1.9$ only under a periodic correction the data reject. Both are robust in sign against population splits, injection--recovery bias correction, hierarchical bootstrap, clustered-injection tests, and an independent FilFinder cross-check. Its precise value carries a factor-$\sim$2 skeleton-algorithm systematic. Direct Plummer fits to all four regions ($W_{\rm fil}=0.11$--$0.16$\,pc, $p\simeq2.0$) give $\lambda/W=1.1$--$1.8$ (median $1.4$), all sub-classical, with Taurus (a classical outlier only under a narrow Gaussian-crest width) the most sub-classical at $\approx1.1$; the corresponding \emph{fixed-width}, periodic-corrected per-region values are $1.0$, $1.9$, $2.1$, $1.8$ (Taurus, Orion~B, Perseus, Aquila; the region-Plummer analogues are lower), and dropping Taurus leaves the median at $1.9$. The pairwise median converges to $L/3$, not the fragmentation wavelength. The core distribution is random-to-clustered rather than periodic (Section~\ref{sec:clustering}), which both fixes the injection--recovery correction to the clustered branch and independently disfavours a single quasi-periodic fragmentation wavelength.
        \item The measurement pipeline is unbiased: applied to a genuinely force-balanced (relaxed, pressure-confined) equilibrium cylinder it recovers the classical $\lambda_{\rm m}=22H\approx5$--$6\,W_{\rm FWHM}$ ($\lambda_{\rm max}/22H=1.07$--$1.10$, converged, matching the analytic Ostriker $22H/W_{\rm FWHM}=5.07$ to $\approx1\%$), whereas the \emph{same} pipeline gives $\approx1.4\,W_{\rm FWHM}$ for the contracting filaments --- so the sub-classical scale is a physical consequence of contraction, not a measurement or normalisation artefact (a closed-loop synthetic test independently recovers injected $\lambda/W$ up to $5.7$ to $<0.5\%$; this does not test the skeletonisation, the dominant observational systematic). The contracting-filament wavelength is sub-classical, $\approx0.35$--$0.5$ times the classical value (a factor $\approx2$--$3$; $\approx1$ background Jeans length): it is an order-$\lambda_{\rm J}$ \emph{upper limit} (the spectrum does not turn over within the converged band; Section~\ref{sec:validation_endtoend}), not a precisely-located maximum. Boundary condition, profile and seed each move it by $\lesssim40\%$.
        \item This is the wavelength expected for a filament that fragments while still contracting, rather than from hydrostatic equilibrium \citep{Clarke2016}. \textbf{The primary conclusion is under-determination, not sufficiency.} What is robust is the \emph{sign}: sub-classical by a factor $\approx3$--$4$ (DisPerSE conventions; more under FilFinder) relative to the classical $\lambda_{\rm m}\approx5$--$6\times$FWHM. The precise value spans a factor of several (Table~\ref{tab:headline}), so whether single-filament fragmentation is by itself \emph{sufficient} depends on the segmentation convention: in matched (beam-convolved Plummer) units the simulated wavelength ($\approx1.0$--$1.3$, an \emph{upper limit}) \emph{matches} under the region-Plummer widths, \emph{under}-predicts under fixed-width DisPerSE ($\sim1.3$ vs $1.9$), and \emph{over}-predicts under FilFinder ($0.4$--$0.8$); because it is an upper limit, a shorter true wavelength would make it over-produce cores. Hierarchical fibres and turbulence remain viable; of the magnetic geometries the static-perpendicular route is disposed of (a resolved perpendicular field spindles rather than beads; Section~\ref{sec:optional_tests}) and the helical/toroidal route --- the only one that could shorten $\lambda$ --- remains \emph{untested}. The discriminating test is a fibre-resolved core spacing across the robust regions.
    \end{itemize}

    \item \textbf{Viable (but not required) explanations.}
    \begin{itemize}
        \item \emph{Hierarchical-fibre projection}: if filaments are bundles of $N$ velocity-coherent fibres of width $W_{\rm fib}$, each fragmenting at its own classical wavelength, the filament-level spacing is $\lambda/W_{\rm fil}\approx(5$--$6)r/N$ ($r=W_{\rm fib}/W_{\rm fil}$; fixed-width DisPerSE requires $r/N\approx0.35$, region-Plummer $\approx0.25$) --- a natural reading of the fibre-resolved data \citep{Yang2024,Hacar2013} that also accounts for the skeleton-algorithm and $L/3$ systematics. But if the fibres themselves fragment sub-classically (as our own result implies for near-critical cylinders) the required $r/N$ roughly doubles and the fibre and single-filament pictures are then \emph{not distinct mechanisms} (Section~\ref{sec:fibre_projection}). Taken at face value the observed fibre counts ($N\approx3$--$5$, $r\approx0.3$--$0.5$, so $r/N\approx0.07$--$0.15$) sit \emph{below} the $r/N\approx0.35$ the DisPerSE spacing requires, so projection \emph{over}-compresses and is admissible only in a narrow low-$N$ corner. Decisively testable with fibre-resolved NN spacing.
        \item \emph{Magnetic tension}: a shorter spacing requires a \emph{helical}/toroidal field \citep{Fiege2000} (static longitudinal or perpendicular fields do not shorten the axial mode; Section~\ref{sec:magnetic}). We partly address this: a resolved perpendicular field spindles rather than beads (disposing of the static-perpendicular route), and a \emph{seeded} helical-field grid ($B_\phi/B_z=0$--$2$) does \emph{not} shorten $\lambda_{\rm max}$ --- but it narrows $W$ by a radial pinch, moving $\lambda/W$ towards (not away from) the classical value, and being out of force balance it does not contain the \citet{Fiege2000} pinch mode. The helical/toroidal mechanism therefore remains \emph{untested}; a force-balanced helical equilibrium (Section~\ref{sec:optional_tests}) is required to settle it and is left to future work.
        \item A resolved perpendicular-field filament ($N_J$ to $\approx170$; Section~\ref{sec:optional_tests}) collapses monolithically to a spindle with no axial fragmentation, consistent with static linear theory (a perpendicular field lengthens the axial mode); the short beading seen at low resolution is a Truelove artefact. Perpendicular filaments therefore do not fragment axially, so they contribute spindles rather than the observed core chains.
    \end{itemize}

    \item \textbf{Speculative.}
    \begin{itemize}
        \item A single illustrative simulation indicates the isothermal assumption can be non-conservative locally under a $\sim12\%$ temperature gradient; a fuller analysis is deferred to a companion paper.
    \end{itemize}

    \item \textbf{Outlook.} The decisive test that would distinguish single-filament from hierarchical-fibre fragmentation is fibre-resolved NN spacing on the full HGBS catalogues; complementary tests are high-resolution polarimetry of filament interiors, a resolution-convergence study of the perpendicular-field case, and MHD runs with a spatially varying temperature field.
    \end{itemize}

\section*{ACKNOWLEDGMENTS}

We thank the Herschel Gould Belt Survey team for the datasets. HGBS data are available from the Herschel Science Archive. The Athena++ simulations were run on cloud computing infrastructure (Google Cloud E2); the full set of $2251$ Athena++ MHD simulations (Appendix~\ref{sec:campaigns}) consumed of order $10^5$ CPU-hours.

\section*{DATA AVAILABILITY}

A public Zenodo archive accompanies this paper (DOI reserved; released with acceptance) and contains: (i) the full run manifest (all 2251 runs with their $(f,\beta,\mathcal{M},\theta,\eta_{\rm AD})$ parameters and outcomes); (ii) the Athena++ problem-generator source (\texttt{filament\_ambient.cpp}, including the divergence-free helical construction) and initial-condition normalisation scripts; (iii) the HDF5 snapshots of the growth-rate convergence subset ($128^3$--$512^3$) and the drag-relaxed equilibrium-recovery control (remaining snapshots, $\sim$0.1\,TB, on request); and (iv) the analysis pipeline --- DisPerSE parameters (v0.9.24; $3\sigma$ persistence, smoothing as in Section~\ref{sec:methodology}), the bridging/deblending and nearest-neighbour scripts, and the derived $\lambda/W$ catalogues. The estimator was independently re-implemented and checked against the deposited counts and medians (Table~\ref{tab:nnpipeline}). HGBS column-density and skeleton maps are available from the Herschel Science Archive.

\bibliographystyle{mnras}
\appendix

\section{Validation of the measurement and the width normalisation}
\label{sec:val_appendix}

This appendix documents the three checks summarised in Section~\ref{sec:validation_endtoend}.

To verify that the routines extracting $\lambda$ and $W$ carry no hidden normalisation bias, we built synthetic filament cubes at the simulation resolution with a \emph{known} transverse FWHM $W_{\rm in}$ and a \emph{known} axial bead spacing $\lambda_{\rm in}$ (localised Gaussian beads on a Gaussian transverse profile) and passed them through the identical pipeline used on the simulations (spine-averaged transverse profile, Gaussian and beam-convolved Plummer width fits, axial peak-spacing). The recovered $\lambda/W$ matches the input to better than $0.5\%$ over $\lambda/W_{\rm in}=1.4$--$5.7$ (Fig.~\ref{fig:pipeval}). The measurement is therefore unbiased, and the corrected width bookkeeping of Eq.~(\ref{eq:t1_apply}) is confirmed.

\begin{figure}
\includegraphics[width=0.46\textwidth]{figures/fig_pipeline_validation.pdf}
\caption{Closed-loop validation: the pipeline recovers the injected $\lambda/W$ to within $0.4\%$ over $\lambda/W=1.4$--$5.7$.}
\label{fig:pipeval}
\end{figure}

\emph{Growth-rate measurement and convergence.} We take the fragmentation wavelength as the peak of the growth-rate spectrum (numerically convergeable), not a late-time peak count (which sub-fragments with resolution). Seeding eight ambient-confined filaments ($f=1.1,1.2,1.5$; $\beta=1$; two seeds) with a broadband perturbation ($\delta v/c_s=10^{-3}$, twelve axial modes) and fitting each mode's growth over $\delta_{\rm rms}=3\times$seed$\to0.3$ gives $\Gamma(\lambda)$ (Fig.~\ref{fig:growthcurves}). The spectrum rises towards shorter $\lambda$ and is resolution-converged for $\lambda\gtrsim1\,\lambda_{\rm J}$: for $f=1.1$ the $n=7$ mode ($\lambda=1.14\,\lambda_{\rm J}$) has $\Gamma=17.8,18.5,18.7$ at $128^3,256^3,512^3$ (stable to $\sim5\%$), and it is flat through the converged-band edge ($\Gamma(n{=}6,7,8)=17.6,18.7,18.7$ at $512^3$: a \emph{plateau, not a turnover}), while modes at $\lambda\lesssim0.9\,\lambda_{\rm J}$ ($n\geq9$) grow with transverse resolution and are numerical. The converged edge does not migrate with resolution, so the $\lambda_{\rm max}\lesssim1\,\lambda_{\rm J}$ \emph{upper bound} is resolution-stable, not resolution-set, and an independent seed reproduces the selected wavelength to $3\%$. Across the eight runs $\lambda_{\rm max}=0.9$--$1.1\,\lambda_{\rm J}$ ($\lambda_{\rm max}/W_{\rm FWHM}=1.5$--$2.0$; Table~\ref{tab:wave}), so $\lambda_{\rm max}/\lambda_{\rm IM92}=0.4$--$0.5$; extending to $f=1.8,2.0$ gives the same sub-classical scale, covering the HGBS line-mass range ($f\approx0.7$--$1.9$; Appendix~\ref{sec:supplementary_campaigns}). Periodic and reflecting boundaries are indistinguishable (KS $p=0.17$), so the wavelength is not a boundary artefact.

\emph{Comparison with linear theory, and the gravity solver.} Figure~\ref{fig:dispersion} overplots the \emph{from-scratch--solved} linear dispersion relation of the isothermal self-gravitating equilibrium cylinder (Ostriker profile, $m=0$; reproducing \citealt{Nagasawa1987,Inutsuka1992,Fischera2012} to $<1\%$: peak at $\lambda_{\rm max}=22.0\,H$, thermal cutoff $\lambda_{\rm cut}=11.1\,H$). In our background-Jeans normalisation $\lambda_{\rm IM92}=22\,H=2.3\,\lambda_{\rm J}$, so the equilibrium spectrum peaks at $2.3\,\lambda_{\rm J}$ and cuts off at $\approx1.16\,\lambda_{\rm J}$. The measured spectrum of the collapsing filament instead rises monotonically and passes \emph{through} the equilibrium cutoff without turning over in the converged band ($\lambda\gtrsim1\,\lambda_{\rm J}$); the small-scale rise ($\lambda\lesssim0.9\,\lambda_{\rm J}$) is identified as numerical because it \emph{diverges with resolution}, not because it lies beyond the cutoff. The dominant physical wavelength is therefore an order-$\lambda_{\rm J}$ \emph{upper limit}, $\lambda_{\rm max}/W_{\rm FWHM}\approx2$ versus the classical $5$--$6$ (\citealt{Fischera2012}), a factor $\approx2.5$--$3$ --- the direction expected for a contracting, non-equilibrium filament \citep{Clarke2016,Clarke2017}. The comparison is made as the density-independent ratio $\lambda/W_{\rm FWHM}$ co-measured at one epoch; self-gravity uses a periodic FFT with reflecting transverse hydrodynamic boundaries, and both a transverse domain-doubling ($2\to4\,\lambda_{\rm J}$) test and an isolated-boundary re-solution of Poisson's equation leave $\lambda_{\rm max}$ unchanged (the box-mean subtraction perturbs only the two longest modes, agreeing to $<1\%$ at $\lambda_{\rm max}$), so the selected wavelength is unbiased by the gravity boundary treatment.

\begin{figure}
\includegraphics[width=0.48\textwidth]{figures/fig_dispersion.pdf}
\caption{The from-scratch--solved eigenvalue dispersion relation of the isothermal self-gravitating equilibrium cylinder (\citealt{Nagasawa1987}; \citealt{Inutsuka1992}), validated to $<1\%$ against \citet{Nagasawa1987} and \citet{Fischera2012}, not a fit: it peaks at the established most-unstable wavelength $\lambda_{\rm IM92}\approx2.3\,\lambda_{\rm J}$ and vanishes at the thermal cutoff $\lambda_{\rm crit}\approx1.2\,\lambda_{\rm J}$, overplotted on the measured growth-rate spectrum ($f=1.1$, three resolutions), each normalised to its value at $\lambda=1.14\,\lambda_{\rm J}$. The analytic curve peaks at the classical value and vanishes below $\lambda_{\rm crit}$; the measured spectrum instead rises monotonically into the grid-scale region (shaded), confirming that the small-scale rise is numerical. The measured wavelength ($\approx1.1\,\lambda_{\rm J}$, dotted; $\lambda_{\rm max}/W_{\rm FWHM}\approx2$) lies a factor $\approx2.5$--$3$ below the analytic equilibrium peak ($\lambda_{\rm m}\approx5$--$6\times$FWHM), the dynamical (contracting-filament) shift, and near the analytic cutoff; we therefore treat it as an order-$\lambda_{\rm J}$, sub-classical value rather than a precisely-located maximum (text).}
\label{fig:dispersion}
\end{figure}

\begin{figure}
\centering
\includegraphics[width=0.99\columnwidth]{figures/fig_eqrecovery.png}
\caption{Equilibrium-recovery control (the primary validation that the sub-classical scale is physical). The identical growth-rate \emph{measurement} pipeline is applied to a \emph{relaxed, force-balanced} pressure-confined Ostriker equilibrium (blue; prepared by explicit drag relaxation, then evolved drag-free) and to a \emph{dynamically contracting} supercritical filament (orange); the two differ in their prepared initial state, not in the measurement, isolating the effect of starting from equilibrium versus contraction. \emph{Left:} onset fragmentation power spectrum; the equilibrium peaks at the classical $\lambda_{\rm m}\approx22H$ ($\approx5.4\,W_{\rm FWHM}$) while the contracting filament peaks sub-classically ($\approx1.4\,W_{\rm FWHM}$, upper limit). \emph{Right:} the dispersion relation $\Gamma(\lambda)$ versus $\lambda/22H$; the equilibrium shows the classical peak-and-decline with its maximum at $22H$, whereas the contracting filament's $\Gamma$ rises monotonically. The pipeline therefore recovers the classical wavelength when the physics is classical, confirming the sub-classical scale of contracting filaments is physical.}
\label{fig:eqrecovery}
\end{figure}

\emph{Nonlinear bead spacing and timescale separation.} Because the growth rate is fitted while the filament is itself contracting, we checked that the linear peak is the scale at which discrete cores actually form. Evolved into the non-linear regime, the near-critical runs develop discrete axial beads whose spacing, $\lambda_{\rm bead}\approx1.1$--$1.6\,\lambda_{\rm J}$ (measured before bead merging at high contrast reduces the count), overlaps the growth-rate range $\lambda_{\rm max}=0.89$--$1.14\,\lambda_{\rm J}$ (Table~\ref{tab:wave}): the two independent measurements agree at the order-$\lambda_{\rm J}$ level, though not to better than the $\sim30\%$ epoch- and seed-dependent scatter of the bead count. The growth is fast: over the fit window the mode e-folds $\sim3$--$4$ times while the central contrast rises from $C\approx20$ to $\gtrsim100$, so the background evolves on a timescale only a few times longer than the mode growth ($\tau_{\rm growth}/\tau_{\rm bg}\approx0.3$--$0.8$), marginal timescale separation as expected for a collapsing rather than a quasi-static filament. Because of this we tested directly whether the extracted spectrum depends on the fit window: restricting the fit to the \emph{early}, near-static phase ($\delta_{\rm rms}<0.05$, where $N_J\gtrsim5$ and the background is closest to hydrostatic) reproduces the full-window growth-rate spectrum to better than $1\%$ at every physical mode, and the peak mode is unchanged, at $f=1.1$ and $1.5$. The measured wavelength is therefore set at the low-contrast start of the window and is not an artefact of fitting the contracting late-time background; $\Gamma(\lambda)$ is the effective growth rate of the collapsing configuration, but the \emph{selected wavelength} is fit-window invariant, and is independently corroborated by the nonlinear bead spacing above.

\emph{Pipeline benchmark against the analytic Jeans rate.} To verify that the growth-rate machinery returns physical rates, we seeded a near-uniform self-gravitating box ($\rho\simeq\rho_0$) and measured $\Gamma$ per axial mode. For the well-resolved long-wavelength modes the measured rates match the analytic uniform-medium value $\Gamma(k)=\sqrt{4\pi G\rho_0-c_s^2k^2}$ to $1$--$7\%$ ($n{=}1$: $6.16$ vs $6.23$; $n{=}2$: $6.24$ vs $6.08$; $n{=}4$: $5.80$ vs $5.44$), confirming the pipeline. This also explains the apparently large filament rates ($\Gamma\approx18$--$26$): $\Gamma(\rho_0)=\sqrt{4\pi G\rho_0}\approx6.3$, so the filament value corresponds to $\rho/\rho_0\approx(18/6.3)^2\approx8$, i.e.\ $\Gamma$ is the \emph{local} free-fall rate at the elevated central density during the fit window, not an anomalous super-free-fall rate.

\emph{Robustness to line-mass, initial profile and field strength.} The order-$\lambda_{\rm J}$ result is not specific to one configuration. Repeating the growth-rate measurement (i) across $\beta=0.3$, $0.5$, $1.0$ (longitudinal, $f=1.1$) gives $\lambda_{\rm max}=1.00$, $1.14$, $0.80\,\lambda_{\rm J}$ --- order $\lambda_{\rm J}$ with only a weak, non-monotonic $\beta$-dependence within the mode-quantisation step ($\Delta\lambda=L_x/[n(n{+}1)]\approx0.13$--$0.15\,\lambda_{\rm J}$ near $\lambda\sim1$), i.e.\ not a resolved physical trend, so the $\beta=1$ headline is representative and the stronger $\beta$-trend of the late-time bead count in Appendix~\ref{sec:supplementary_campaigns} (which we do not adopt for the wavelength) is a late-epoch artefact; and (ii) replacing the Plummer-like ($p\simeq2$) initial profile with a Ostriker equilibrium profile ($p=4$) gives $\lambda_{\rm max}=0.80\,\lambda_{\rm J}$ at $128^3$ and $256^3$ --- again order $\lambda_{\rm J}$, so the sub-classical scale is not an artefact of the adopted profile index. Notably even the $p=4$ cylinder fragments at $\approx0.8\,\lambda_{\rm J}$ (an order-$\lambda_{\rm J}$, sub-classical scale), not near the classical $\lambda_{\rm m}\approx5$--$6\times$FWHM. We caution, however, that this $p=4$ run is only \emph{approximately} in equilibrium in our discretised setup: its central contrast relaxes by a factor $\sim10$ before beading, so it too fragments from a contracted, dynamical state. We therefore cannot cleanly separate a genuine contraction effect from this initial relaxation, and we report the sub-classical scale as a generic outcome of fragmentation from a non-equilibrium, contracting state --- \emph{not} as a demonstration that a true equilibrium cylinder cannot fragment at the classical scale (which would require a properly relaxed, force-balanced equilibrium control we do not achieve here; Section~\ref{sec:limitations}).

\emph{Recovery of the classical limit for a genuine equilibrium --- the decisive control.} The critical test is whether the \emph{identical} pipeline recovers the classical $\lambda_{\rm m}=22H\approx5$--$6\,W_{\rm FWHM}$ for a genuinely force-balanced equilibrium cylinder. It does. Because the periodic FFT solver enforces the Jeans swindle at the box-mean density, the analytic Ostriker profile is not itself the solver's equilibrium (initialised directly it breathes by factors of two), so we relaxed a pressure-confined Ostriker profile to the solver's hydrostatic state with an explicit velocity drag, verified with a no-perturbation control (central density drifts only $\sim3\%$, in the $k=0$ mode --- not the $C\sim10^4$ runaway of the contracting runs), and seeded it with the identical 12-mode perturbation (box $L_x=16\,\lambda_{\rm J}$, two resolutions). Its $\Gamma(\lambda)$ shows the textbook \emph{peak-and-decline} with maximum at $\lambda_{\rm max}\approx22H$: $\lambda_{\rm max}/W_{\rm FWHM}=5.0$--$5.5$ and $\lambda_{\rm max}/22H=1.07$--$1.10$ (converged; the profile invariant $22H/W_{\rm FWHM}=4.96$--$5.03$ matches the analytic Ostriker $5.07$ to $\approx1\%$). The $\approx7$--$10\%$ long peak is a conservative bias (it would make the contracting deficit appear smaller). Applying the \emph{same} pipeline to the contracting filaments instead gives $\lambda_{\rm max}\approx1\,\lambda_{\rm J}\approx1.4\,W_{\rm FWHM}$ ($\approx0.4\times22H$; the near-critical $f=1.1$ case gives $\approx2$ in $W_{\rm FWHM}$ units --- both sub-classical) with a monotonically rising $\Gamma$ (Fig.~\ref{fig:eqrecovery}). The two experiments differ in their \emph{prepared initial state} (a drag-relaxed force-balanced equilibrium versus a contracting profile) but use the \emph{identical}, drag-free growth-rate measurement; the drag acts only in preparing the equilibrium, not during the measured evolution. On this controlled comparison the pipeline is demonstrably unbiased and the sub-classical spacing is a physical consequence of starting from contraction rather than equilibrium, not a measurement artefact; the closed-loop synthetic test independently recovers injected $\lambda/W$ up to $5.7$ to $<0.5\%$. The control validates the pipeline and the \emph{sign} of the deficit; the contracting $\lambda_{\rm max}$ itself remains an upper limit (its ``peak'' is the converged-band edge, not a resolved maximum).

\begin{table}
\caption{Growth-rate (dominant-mode) measurements. $\lambda_{\rm max}$ is the peak of $\Gamma(\lambda)$, $\Gamma_{\rm peak}$ its growth rate, $W$ the \emph{per-run} Gaussian formation FWHM ($\approx0.55$--$0.67\,\lambda_{\rm J}$, not a single fixed value; all in $\lambda_{\rm J}$). The $f=1.8$, $2.0$ rows extend the \emph{same} growth-rate method into the supercritical regime (single resolution, so not part of the convergence ladder); they give the same sub-classical scale as the bead-count supercritical ensemble ($\lambda/\lambda_{\rm J}\approx0.98$; Appendix~\ref{sec:supplementary_campaigns}), so the two wavelength-measurement methods agree at the $\sim\lambda_{\rm J}$ level where they overlap. The wavelength resolution is set by the discrete axial modes ($\lambda=L_x/n$, $\Delta\lambda\approx0.15\,\lambda_{\rm J}$ near $\lambda\sim1$), which bounds the precision of $\lambda_{\rm max}$. Because the axial modes are discrete ($\lambda=L_x/n$), $\lambda_{\rm max}$ carries a quantisation uncertainty of $\pm\tfrac12\Delta\lambda\approx\pm0.07\,\lambda_{\rm J}$ near $\lambda\sim1$; the final quoted digit is not significant, and the box cannot be lengthened to refine it because $L_x\geq24\,\lambda_{\rm J}$ introduces a seeding artefact (Appendix~\ref{sec:supplementary_campaigns}).}
\label{tab:wave}
\footnotesize
\setlength{\tabcolsep}{4pt}
\begin{tabular}{cccccc}
\toprule
$f$ & resolution & seed & $\lambda_{\rm max}$ & $\Gamma_{\rm peak}$ & $\lambda_{\rm max}/W$ \\
\midrule
1.1 & $128^3$ & 42  & 1.14 & 17.8 & 1.95 \\
1.1 & $256^3$ & 42  & 1.14 & 18.5 & 1.99 \\
1.1 & $512^3$ & 42  & 1.14 & 18.7 & 2.00 \\
1.2 & $256^3$ & 42  & 1.00 & 19.1 & 1.81 \\
1.2 & $256^3$ & 137 & 0.89 & 20.1 & 1.60 \\
1.5 & $256^3$ & 42  & 0.89 & 21.2 & 1.47 \\
1.5 & $256^3$ & 137 & 0.89 & 22.1 & 1.47 \\
1.5 & $512^3$ & 42  & 0.89 & 21.5 & 1.47 \\
1.8 & $256^3$ & 42  & 0.80 & 24.6 & 1.19 \\
2.0 & $256^3$ & 42  & 0.80 & 26.3 & 1.20 \\
\bottomrule
\end{tabular}
\end{table}

A complementary end-to-end grid of 21 ambient-confined runs ($f=1.5$--$3.0$, $\beta=0.5$--$2.0$, two seeds; the forward-modelled subset of the ambient-confined supercritical ensemble, campaigns~11--13; Table~\ref{tab:campaigns}), forward-modelled through beam convolution and Plummer fitting, gives a consistent picture: the spacing in background-density Jeans lengths is $\lambda/\lambda_{\rm J}(\rho_0)\approx0.9$--$1.2$, and the same growth-rate/late-peak distinction accounts for its scatter. Boundary condition, profile and seed each move the value by $\lesssim40\%$ (medians $1.2$ reflecting, $0.8$ periodic, $1.4$ user; $1.4$ Gaussian, $0.9$ Ostriker), and the central-to-ambient contrast varying by a factor $\approx1.6$ across the grid shifts it by only $\sim0.3$, so the wavelength is a property of the filament, not of the imposed confinement. All per-run measurements (configuration, resolution, $\lambda_{\rm max}$, growth rate, epoch) are released with the public data archive.

\emph{Numerical validation (summary).} The fragmentation \emph{classification} is resolution-independent: six parameter points run at both $128^3$ and $256^3$ fragment at both resolutions, with a mean timescale ratio $t_{\rm frag}(256^3)/t_{\rm frag}(128^3)=0.915\pm0.012$ (an $8.5\%$ offset; the collapse \emph{time} itself is not claimed converged), and replacing the Gaussian initial condition with uniform density gives identical classifications. The local Jeans length stays resolved at the epoch the spacing is measured for the headline longitudinal runs ($N_J\approx5$--$7$; \citet{Truelove1997} criterion $N_J\gtrsim4$ satisfied), and the resolved ideal-perpendicular refinement ladder reaches $N_J\approx40$--$170$ (Section~\ref{sec:optional_tests}); only the perpendicular \emph{ambipolar} runs remain unconverged ($N_J\approx1$--$2$), which is why their short wavelengths are treated qualitatively. The full resolution, Truelove/$N_J$ and equation-of-state convergence tables, and the complete 2251-run campaign list, are given in Appendix~\ref{sec:validation} and Table~\ref{tab:campaigns}.


% ===== Reintegrated supporting appendices =====
\section{Simulation Validation Details}
\label{sec:validation}

To ensure robustness against numerical effects we performed three validation tests (83 simulations). For resolution, six parameter points run at both $128^3$ and $256^3$ with matched problem-generator settings, initial conditions, seeds and boundaries all fragment at both resolutions; the mean timescale ratio is
\begin{equation}
\frac{t_{\rm frag}(256^3)}{t_{\rm frag}(128^3)} = 0.915 \pm 0.012,
\end{equation}
an $8.5\%\pm1.2\%$ offset (from the ratio $0.915\pm0.012$), all points within a $\pm11\%$ band (Fig.~\ref{fig:resolution}), so the fragmentation classification is resolution-independent. Replacing the Gaussian-profile initial condition with uniform density yields identical classifications at all 10 tested points.

A fragmentation-wavelength measurement is trustworthy only if the local Jeans length remains resolved ($\gtrsim4$ cells; Truelove et al.\ 1997) up to the epoch at which the spacing is measured. The Truelove criterion was derived to prevent \emph{artificial} fragmentation in self-gravitating collapse; we apply it here as a necessary resolution-adequacy condition for a growth-rate measurement, because an under-resolved Jeans length spuriously enhances small-scale power, exactly the grid-scale, resolution-divergent rise we identify as numerical (Fig.~\ref{fig:growthrate}a). We therefore tracked the minimum cells-per-local-Jeans-length, $N_J$, along each collapse trajectory (Table~\ref{tab:truelove}). For the \emph{longitudinal} runs (both the near-critical calibration and the supercritical ensemble that provides the headline longitudinal value), the inter-bead wavelength ($\lambda\approx\lambda_{\rm J}$, itself resolved by $\sim32$ cells) is established at moderate density contrast ($C\approx20$--$40$ near-critical; early beading well before runaway for the supercritical ensemble), where $N_J\approx5$--$7$ (Truelove satisfied); $N_J$ drops below 4 only at $C>170$, i.e. \emph{after} the spacing is first set. For the perpendicular \emph{ambipolar} runs, by contrast, beading is measured only at $C\approx250$--$560$ where $N_J\approx1$--$2$ (Truelove violated), and the short perpendicular values are \emph{resolution-limited upper limits}. This is why we base the quoted wavelength on the linear growth-rate spectrum rather than on late-time peak counts: the growth-rate peak is established when $N_J\gtrsim5$ and is resolution-converged (the $128^3$--$512^3$ ladder above), whereas the late-time peak count is set by sub-fragmentation at $N_J<4$ and is not.

\begin{table}
\caption{Jeans (Truelove) resolution: minimum cells per local Jeans length $N_J$ at the epoch the spacing is measured, and the density contrast $C$ at that epoch.}
\label{tab:truelove}
\footnotesize
\setlength{\tabcolsep}{4pt}
\begin{tabular}{p{3.6cm}cc}
\toprule
Configuration & $C$ at measurement & min $N_J$ \\
\midrule
Near-critical longitudinal (calibration) & $20$--$40$ & $5$--$7$ (OK) \\
Supercritical longitudinal ensemble (headline) & early beading, moderate $C$ & $\gtrsim4$ (OK) \\
Ideal perp., converged ladder ($512^3$; Section~\ref{sec:optional_tests}) & $\approx3$ / $\approx30$ & $170$ / $40$ (OK; monolithic) \\
Ideal perp., short-$\lambda$ artefact onset & $>10^2$ (grid-dependent) & $\lesssim1$ (violated) \\
Perp.\ ambipolar (tension side) & $250$--$560$ & $1$--$2$ (violated) \\
\bottomrule
\end{tabular}
\end{table}

\begin{figure}
\includegraphics[width=0.48\textwidth]{figures/figR1_resolution_scatter.pdf}
\caption{Resolution convergence with matched problem generator: $t_{\rm frag}(256^3)$ versus $t_{\rm frag}(128^3)$ for six parameter points, all within a $\pm11\%$ band (mean ratio $0.915\pm0.012$).}
\label{fig:resolution}
\end{figure}

\subsubsection{Equation of State Sensitivity}

Testing mildly sub-isothermal equations of state with $\gamma = 0.9$ and $0.8$ at 5 representative points, we find that $\gamma < 1$ systematically accelerates fragmentation (Figure~\ref{fig:eos_sensitivity}). At $\gamma = 0.8$, all 5 test points fragmented at $t_{\rm frag} \approx 0.66$--$0.68\,t_{\rm J}$, approximately half the isothermal fragmentation timescale for the same $(f, \beta, \mathcal{M})$; at $\gamma = 0.9$ fragmentation is delayed (1 of 5 fragments by the integration limit). These are relative comparisons at matched $t_{\rm J}$: the point is that $\gamma<1$ shortens $t_{\rm frag}$ monotonically, consistent with the isothermal runs of the main text (which fragment at $t_{\rm frag}\approx1.1$--$1.2\,t_{\rm J}$); we draw no stability inference from the finite integration window.

For $\gamma<1$ the pressure response to compression weakens ($P\propto\rho^{\gamma}$ rises more slowly than isothermal), reducing Jeans support and hastening fragmentation. Far-IR-cooled cloud filaments typically have $\gamma\approx0.7$--$0.95$, so the isothermal predictions are conservative: real filaments are if anything more fragmentation-prone.

\begin{figure}
\includegraphics[width=0.48\textwidth]{figures/fig6_phase3_eos_sensitivity.pdf}
\caption{Equation of state sensitivity test. Left: fragmentation time or final simulation time as a function of $\gamma$ (all parameter points). Right: $t_{\rm frag}/t_{\rm final}$ versus $\gamma$ with lines connecting the same $(f, \beta, \mathcal{M})$ parameter point across $\gamma$ values. Sub-isothermal EOS ($\gamma < 1$) systematically reduces $t_{\rm frag}$.}
\label{fig:eos_sensitivity}
\end{figure}

\section{Complete simulation campaign list}
\label{sec:campaigns}

Table~\ref{tab:campaigns} lists every simulation campaign cited in the manuscript. Rows 5--8 are the components of what was previously reported as a single Field-Geometry campaign (314 runs); they are counted individually here and the 314 aggregate is \emph{not} added separately.

\begin{table*}
\caption{Complete list of simulation campaigns underlying the paper.}
\label{tab:campaigns}
\begin{tabular}{cllrl}
\toprule
\# & Campaign (section) & Coverage $(f,\beta,\mathcal{M},\theta,\eta_{\rm AD})$ & Runs & Independent? \\
\midrule
1  & Base three-regime grid (\S4.5; Fig.~\ref{fig:regime}) & $f\approx1$; $\beta\,0.05$--$5$; $\mathcal{M}\,0.5$--$5$; $\theta=0$; ideal & 208 & yes \\
2  & Transition grid / DTC & $f\,1.4$--$2.2$; $\beta\,0.3$--$1.3$; $\mathcal{M}\,1$--$5$; $\theta=0$ & 540 & yes \\
3  & Supercritical grid (\S4.6) & $f\,1.1$--$3.0$; $\beta\,0.3$--$5.0$; $\mathcal{M}\,0.5$--$3$; periodic & 654 & yes \\
4  & Validation (\S4.4; App.~\ref{sec:validation}) & resolution / IC / EOS convergence & 83 & yes \\
5  & Near-critical (\S4.7.1) & $f\,1.0$--$1.2$; $\beta\,0.3$--$1.0$; $\mathcal{M}\,1$--$2$; $\theta=0$ & 80 & FG component \\
6  & Perpendicular (\S4.7.2) & $\theta=90^\circ$; $f\,1.2$--$1.5$; $\beta\,0.3$--$2.0$ & 96 & FG component \\
7  & Oblique calibration (\S4.7.3) & $\theta=30/45/60^\circ$; $\beta\,0.5$--$2.0$ & 108 & FG component \\
8  & Adiabatic EOS & $\gamma=5/3$; $f\,1.0$--$1.2$ & 30 & FG component \\
9  & Turbulence + perp-$\beta$ + crit.-transition (\S4.8; incl.\ turbulent-support $f_{\rm eff}$ 90 + physical-vs-synthetic turbulence comparison 54) & $\delta v/c_s\,10^{-4}$--$1$; $\beta\,0.3$--$2$ & 289 & yes (aggregate) \\
10 & Targeted DTC re-runs & $\beta=0.3,\mathcal{M}=1$; extended timeout & 25 & yes \\
11 & Direct supercritical ensemble & $f\,1.5$--$2.0$; $\beta\,0.5$--$2.0$; $\theta=0$; ambient & 39 & yes \\
12 & Reconciliation suite (\S4.6.6) & $f\times\beta\times2$ seeds; $\theta=0$; ambient & 12 & yes \\
13 & Ideal-MHD perp.\ $d=1.0$ grid (\S4.6.6) & $\theta=90^\circ$; $\beta\,0.5$--$2.0$; ambient & 18 & yes \\
14 & Ambipolar non-ideal survey & $\theta=90^\circ$; $\eta_{\rm AD}\,0.001$--$0.1$; $\beta\,0.5$--$2.0$ & 36 & yes \\
15 & Growth-rate + equilibrium control + transverse/higher-$f$ tests (\S\ref{sec:validation_endtoend}, App.~\ref{sec:val_appendix}) & $f\,0.7$--$2.0$; $128^3$--$512^3$; Gaussian/Ostriker; transverse $2$--$4\,\lambda_{\rm J}$ & 25 & yes \\
16 & Axial-resolution stability (\S\ref{sec:optional_tests}) & $f=1.1$; $\theta=0$; $1024\times128^2$ growth-rate rerun & 1 & yes \\
17 & Ideal-perpendicular refinement ladder (\S\ref{sec:optional_tests}; Table~\ref{tab:truelove}) & $\theta=90^\circ$; $f=1.1$; $128^3$--$512^3$ & 3 & yes \\
18 & Divergence-free helical grid (\S\ref{sec:optional_tests}) & $f=1.1$; $B_\phi/B_z=0$--$2$; $512\times128^2$ & 4 & yes \\
\midrule
   & \textbf{Athena++ MHD total} & & $\boldsymbol{2251}$ & \\
\bottomrule
\end{tabular}
\end{table*}

\section{Supplementary simulation-campaign results}
\label{sec:supplementary_campaigns_app}

The campaigns in this appendix are exploratory and are \emph{not} used quantitatively in the main conclusions; the headline wavelength rests solely on the near-critical longitudinal growth-rate ladder of Section~\ref{sec:validation_endtoend}.

\subsection{The radial-collapse/fragmentation timescale competition}
\label{sec:regimes}

The simulations show that the fate of an isothermal, self-gravitating filament is governed by the competition between radial collapse and longitudinal fragmentation \citep{Inutsuka1997, Pon2012}. Near-critical filaments ($f\lesssim1.2$) retain sufficient thermal and magnetic support to evolve quasi-statically, and the longitudinal fragmentation instability grows to non-linearity, producing the beading pattern from which $\lambda/W$ can be measured directly. Supercritical filaments ($f\gtrsim1.5$) fragment longitudinally when properly confined by an ambient medium (Section~\ref{sec:modecompetition}), at a converged \emph{upper bound} on the wavelength $\lesssim0.5$ times the classical value (Section~\ref{sec:validation_endtoend}). Without ambient confinement the filament's Gaussian skirt feeds an accretion channel through the boundary and suppresses fragmentation. Extended-integration tests confirm that the periodic grid's early termination fires before the longitudinal instability becomes non-linear ($\sigma_{\rm lon}$ onset at $t=0.42$--$0.71\,t_{\rm J}$ versus a Courant-limited kill at $t\approx0.25$--$0.34\,t_{\rm J}$), independent of resolution.

In the near-critical regime ($f\approx1.0$--$1.2$), filaments fragment with clearly measurable longitudinal density peaks at $t_{\rm frag}\approx1.1$--$1.2\,t_{\rm J}$ (Section~\ref{sec:modecompetition}). Runs at the exact Jeans-critical value ($f=1.00$) likewise bead, at $t_{\rm frag}=1.62\pm0.03\,t_{\rm J}$ across seeds and resolutions (this exact-critical, weak-field subset fragments later than the $f=1.0$--$1.2$ field-geometry grid above, whose stronger seeding and field bring it to $1.1$--$1.2\,t_{\rm J}$), so there is no stability threshold at the critical line mass; this is the regime in which $\lambda/W$ is directly measurable.

\textbf{Supercritical regime ($f \geq 1.5$): Radial collapse dominates}. In the periodic supercritical grid ($f=1.5$--$3.0$) we find no longitudinal fragmentation structure, only pure radial collapse: the central density increases by two orders of magnitude while the longitudinal density variance remains at $<0.02\%$ up to and including the fragmentation time. All simulations undergo runaway radial collapse (detected as $\Delta t < 10^{-8}\,t_{\rm J}$ from the CFL limit) before longitudinal beading can develop.

The supercritical radial-collapse time ($t_{\rm coll}\approx0.25$--$0.34\,t_{\rm J}$) is $3$--$5$ times shorter than the near-critical longitudinal-fragmentation time ($t_{\rm frag}\approx1.1$--$1.6\,t_{\rm J}$), so in configurations whose initial conditions are out of radial equilibrium the radial mode reaches runaway before longitudinal perturbations grow to non-linearity. The measurable fragmentation wavelengths therefore come from two regimes: the near-critical, quasi-static regime where longitudinal beading is observed directly (Section~\ref{sec:modecompetition}), and ambient-confined supercritical filaments that fragment longitudinally before radial collapse (Section~\ref{sec:modecompetition}).

\subsection{Three fragmentation regimes}

The baseline grid, at a near-critical line mass ($f\approx1$), reveals three fragmentation regimes in the $(\beta,\mathcal{M})$ plane, characterised by the end-of-run density contrast $C=\rho_{\rm max}/\rho_0$ (Fig.~\ref{fig:regime}): (i) \emph{strong-field, magnetically supported} ($\beta\lesssim0.15$), where magnetic pressure strongly suppresses collapse ($C\approx1$ at the snapshot time); (ii) \emph{magnetically regulated} ($0.2\lesssim\beta\lesssim2$), where fields modify but do not prevent fragmentation ($C=3$--$11$); and (iii) \emph{thermally dominated} ($\beta\gtrsim3$), where fields are too weak to affect the scale ($C=13$--$22$). The dispersion relation \citep{Nakamura1993} predicts near-complete suppression for $\beta\lesssim0.15$ and $<10\%$ modification of the most unstable wavelength for $\beta\gtrsim2$--$3$, consistent with the empirical boundaries; these boundaries bracket the sampled grid points ($\beta=0.05,0.1,0.15,0.2,0.3,\ldots$) at roughly the one-grid-step level rather than being interpolated, so the $\beta\approx0.15$ and $\approx3$ thresholds carry an uncertainty of $\sim$one grid step. Observed dense-filament conditions ($\mathcal{M}\approx1$--$3$, $\beta\approx0.3$--$2$; \citealt{Arzoumanian2019,Planck2016,Pattle2018}) place most HGBS filaments in the magnetically regulated regime.

\begin{figure*}
\includegraphics[width=0.8\textwidth]{figures/fig2_regime_corrected.pdf}
\caption{Three-regime fragmentation framework from the Athena++ base grid (near-critical, $\theta=0^\circ$, isothermal; Section~\ref{sec:regimes}). Colour shows the end-of-run density contrast $C=\rho_{\rm max}/\rho_0$ (blue = low, red = high). The regimes, set by plasma $\beta$, are: (I) \emph{strong-field, magnetically supported} ($\beta\lesssim0.15$): magnetic pressure strongly suppresses collapse, $C\approx1$ (minimal fragmentation at this epoch, not an absence-of-star-formation statement); (II) \emph{magnetically regulated} ($0.2\lesssim\beta\lesssim2$): fields modify but do not prevent fragmentation, $C=3$--$11$; (III) \emph{thermally dominated} ($\beta\gtrsim3$): fields too weak to affect the scale, $C=13$--$22$. Typical HGBS filament conditions ($\mathcal{M}\approx1$--$3$, $\beta\approx0.3$--$2$; \citealt{Arzoumanian2019,Planck2016,Pattle2018}) lie in the magnetically regulated regime. The regime here is set by the \emph{dynamical} suppression of collapse (a $\beta$ criterion) and is physically distinct from the magneto-static mass-to-flux criticality (Section~\ref{sec:regimes}); the two are not in tension.}
\label{fig:regime}
\end{figure*}

\subsection{Turbulence and critical-transition dependence}
\label{sec:supplementary_campaigns}

Additional grids map the turbulence, field-geometry and critical-transition dependence of $\lambda/W$.

\subsubsection{Turbulence versus $\lambda/W$}

Comparing physical turbulence ($\delta v/c_s=1.0$, Kolmogorov) with synthetic perturbations ($10^{-4}$) at longitudinal field ($f=1.0$--$1.2$, $\beta=0.5,1,2$), the fragmentation \textit{wavelength} is unchanged ($<5\%$; $\lambda/W=3.44\pm0.76$ vs $3.30\pm0.85$; Fig.~\ref{fig:C5_turbulence}), with a $\beta$-dependence in the \emph{raw late-time bead-count} ratio ($4.31\pm0.60$ at $\beta=0.5$ to $2.81\pm0.13$ at $\beta=2.0$); as the growth-rate analysis shows (Section~\ref{sec:validation_endtoend}), this trend is a late-epoch bead-count effect, not a shift in the underlying $\lambda_{\rm max}$, whose $\beta$-variation is weak and within the mode-quantisation step. Turbulence does accelerate collapse by $3.2\times$ ($t_{\rm frag}=0.33\pm0.06$ vs $1.06\pm0.16\,t_{\rm J}$). As a single decaying realisation, this constrains, but does not resolve, whether \emph{driven} transonic turbulence is a first-order parameter for the fragmentation scale; we conclude only that the characteristic spacing is insensitive to the seed perturbation amplitude in this limited test, and do not claim turbulence is generically unimportant.

\begin{figure*}
\includegraphics[width=0.8\textwidth]{figures/fig_C5_turbulence_lambdaW.pdf}
\caption{turbulence comparison: physical vs synthetic turbulence give statistically indistinguishable $\lambda/W_{\rm core}$ ($\Delta\mu<5\%$, KS $p=0.43$). The $\beta$-trend in the raw late-time bead-count ratio ($\sim$4.3 at $\beta=0.5$ to $\sim$2.8 at $\beta=2.0$) is preserved across the two turbulence prescriptions, but is a late-epoch bead-count effect rather than a shift in the underlying growth-rate $\lambda_{\rm max}$, whose $\beta$-variation is weak (Section~\ref{sec:validation_endtoend}).}
\label{fig:C5_turbulence}
\end{figure*}

\subsubsection{Perpendicular-field spacing}

For perpendicular fields ($\theta=90^\circ$, $f=1.2$--$1.5$, $\beta=0.3$--$2.0$; extended $16\lambda_{\rm J}$ periodic axial domain): strong perpendicular B ($\beta\leq0.5$) gives pure radial collapse with no axial fragmentation, while for $\beta\geq1.0$ axial beading occurs at $\lambda/W_{\rm core}=1.25\pm0.09$ ($N=27$), i.e.\ $\lambda/W_{\rm fil}\approx0.4$ under the corrected conversion (Eq.~\ref{eq:t1_apply}), a factor $2$--$3$ \textit{shorter} than the longitudinal-field value at the same $\beta$. Field geometry, not $\beta$, dominates, and the perpendicular spacing lies far below the observed NN value (fixed-width, periodic-corrected, $\lambda/W\approx1.9$; raw $2.1$).

The two perpendicular campaigns give different $\beta$-trends: this periodic-domain grid beads only for $\beta\geq1.0$, whereas the ambient-confined ideal-MHD grid of Section~\ref{sec:modecompetition} beads at $\beta=0.5$ and $2.0$ with spindle collapse near $\beta\approx1$. They differ in transverse confinement and are both numerically unresolved ($N_J\approx1$--$2$; Appendix~\ref{sec:validation}), so we draw no physical $\beta$-trend from either; a matched-resolution, matched-confinement grid is needed to settle it (Section~\ref{sec:limitations}). The robust statement is only that the perpendicular beading wavelength, where it occurs, is short ($\lambda/W_{\rm fil}\lesssim0.6$), well below the observed spacing.

\subsubsection{Critical-transition spacing}

Mapping $\lambda/W$ across $f=0.9$--$1.3$ at five $\beta$ values (Fig.~\ref{fig:C7_critical}) shows no discontinuity at $f=1.0$ (variation $<5\%$) and an apparent $\beta$-dependence in the \emph{late-time bead count}: $\lambda/W_{\rm core} = 4.74\pm0.73,\,4.38\pm0.59,\,3.19\pm0.29,\,2.86\pm0.18,\,2.80\pm0.14$ at $\beta=0.3,0.5,1.0,1.5,2.0$. This bead-count $\beta$-trend is, however, larger than the primary growth-rate measurement: repeating the growth-rate-spectrum method (Section~\ref{sec:validation_endtoend}) across $\beta=0.3,0.5,1.0$ gives $\lambda_{\rm max}=1.00,\,1.14,\,0.80\,\lambda_{\rm J}$, i.e.\ order one Jeans length with only a weak, mode-quantised $\beta$-dependence. We therefore treat the $\beta$-dependence of the wavelength as weak (the stronger bead-count trend being a late-epoch, resolution-sensitive artefact) and the two highest-$\beta$ bead-count points ($\lambda/\lambda_{\rm J}\approx0.84$) as consistent with the supercritical ensemble; the observable $\lambda/W_{\rm fil}$ then follows from the end-to-end forward model (Section~\ref{sec:validation_endtoend}) rather than the shortcut factor. \textbf{Continuity across $f\approx1$ does not validate extrapolation to the supercritical regime ($f\geq1.5$)}, where the mode changes to radial collapse (Section~\ref{sec:validation_endtoend}); the near-critical $\lambda/W$ values are a hypothesis for supercritical filaments, not a measurement.

\begin{figure*}
\includegraphics[width=0.8\textwidth]{figures/fig9_nearcritical_corrected.pdf}
\caption{Near-critical runs: $\lambda/W_{\rm core}$ across $f=0.9$--$1.3$ for five $\beta$ values, smooth with no discontinuity at $f=1.0$ (left); mean per $\beta$ (right). The convergence-independent quantity is $\lambda/W_{\rm core}$; the observable follows from $\lambda/W_{\rm fil}\approx0.27\,(\lambda/W_{\rm core})$ (Eq.~\ref{eq:t1_apply}), so the highest-$\beta$ points correspond to $\lambda/W_{\rm fil}\approx0.75$--$0.8$. This near-critical calibration lies below the observed spacing under any convention (Table~\ref{tab:headline}); the primary simulated value is the growth-rate wavelength (an upper limit; Section~\ref{sec:validation_endtoend}), not this ratio.}
\label{fig:C7_critical}
\end{figure*}

\begin{table}
\caption{Effective line-mass fraction with turbulent support.}
\label{tab:feff}
\footnotesize
\setlength{\tabcolsep}{4pt}
\begin{tabular}{p{3.0cm}cc}
\toprule
Quantity & Value & Source \\
\midrule
Nominal $f$ (thermal) & $1.5 \pm 0.3$ & HGBS line-mass \\
Turbulent support $f_{\rm eff}/f$ & $0.82 \pm 0.11$ & TURB (90 sims) \\
Effective $f_{\rm eff}$ & $1.23 \pm 0.28$ & product \\
HGBS ensemble range & $0.7$--$1.9$ & propagated \\
\bottomrule
\end{tabular}
\end{table}

\subsubsection{Synthesis}

Table~\ref{tab:cross_campaign} summarises the simulation $\lambda/W$ measurements and their role in the comparison with observation.

\begin{table*}
\caption{Simulation fragmentation scale by configuration, in the two convention-independent forms: the raw internal ratio $\lambda/W_{\rm core}$ (normalised by the initial half-width) and the background-density Jeans-length ratio $\lambda/\lambda_{\rm J}=0.3\,(\lambda/W_{\rm core})$. We do \emph{not} tabulate the observable $\lambda/W_{\rm fil}$ here, since it depends on the width convention; the comparison with observation uses the growth-rate wavelength expressed in \emph{matched} beam-convolved Plummer units, $\lambda/W\approx1.0$--$1.3$ ($\approx0.35$--$0.5$ times the classical value, a factor $\approx2$--$3$; Section~\ref{sec:validation_endtoend}), which is set against the observed $\lambda/W\approx0.4$ (FilFinder) to $1.9$ (fixed-width DisPerSE).}
\label{tab:cross_campaign}
\footnotesize
\setlength{\tabcolsep}{4pt}
\begin{tabular}{lccl}
\toprule
Configuration & $\lambda/W_{\rm core}$ & $\lambda/\lambda_{\rm J}$ & Status \\
\midrule
\multicolumn{4}{l}{\emph{Longitudinal field (quantitative)}} \\
Long.\ near-critical ($\beta{=}2.0$, calibration) & $2.80{\pm}0.14$ & $0.84$ & Calibration only \\
Long.\ supercritical ensemble (\S\ref{sec:modecompetition}) & $3.27$ & $0.98$ & Near one Jeans length; see growth-rate result \\
\midrule
\multicolumn{4}{l}{\emph{Perpendicular field --- PRELIMINARY, numerically unconverged ($N_J{\approx}1$--$2$); not for quantitative use}} \\
Perp.\ near-critical ($\beta{\geq}1$, ideal-MHD) & $[1.25]$ & $[0.38]$ & unconverged \\
Perp.\ ambipolar ($\eta_{\rm AD}{=}0.001$--$0.1$) & $[{\sim}2.2]$ & $[0.66]$ & unconverged \\
\bottomrule
\end{tabular}
\\[2pt]
Only the longitudinal rows are used quantitatively, and even there the reference value is the convergently-bounded growth-rate wavelength (an upper limit; Section~\ref{sec:validation_endtoend}, Table~\ref{tab:wave}), not these raw ratios. The perpendicular values are bracketed to signal that they are numerically unconverged ($N_J\approx1$--$2$, below the Truelove criterion; wavelengths shift by factors $2$--$5$ between resolutions) and must not be read as measurements.
\end{table*}

With the exception of the 54-run physical-versus-synthetic turbulence \emph{comparison} (Section~\ref{sec:supplementary_campaigns}; distinct from the 90-run turbulent-\emph{support} $f_{\rm eff}$ grid of Table~\ref{tab:feff}, both components of the 289-run aggregate campaign~9 of Table~\ref{tab:campaigns}), the grid uses near-laminar seeding ($\delta v=\mathcal{M}c_s\times10^{-4}$) rather than driven ISM turbulence. That comparison shows the fragmentation \textit{wavelength} is unchanged ($<5\%$) at driven amplitudes, so the laminar-seeded $\lambda/W$ values are reported, though the collapse \textit{timescales} are shorter by a factor of $3$--$5$.

In convergence-independent units the two first-order parameters are field geometry (perpendicular $\lambda/W_{\rm core}\approx1.25$ versus longitudinal $2.8$--$4.4$) and, for longitudinal fields, plasma $\beta$; the observable comparison uses the growth-rate wavelength (Section~\ref{sec:validation_endtoend}), not these raw ratios.

\bibliography{references_complete}

\end{document}
